% ch1.tex -- Module 5, Chapter 1: Projectile Motion
% figure_coords.py and check_displayed.py, plus the editor pass
% Convention in the whole chapter: origin at the launch point, up and forwards
% positive, so a_x = 0 and a_y = -g = -9.8 m s^-2.
\chapter{Projectile Motion}\label{ch1:chap}
\covers{M5.IQ1.01, M5.IQ1.02, M5.IQ1.03, M5.IQ1.04}

\inquiry{How can models that are used to explain projectile motion be used to
analyse and make predictions?}

\begin{goals}
  \item split a projectile's motion into a horizontal part and a vertical part
  \item derive and use the maximum height, time of flight, range and final velocity, for a
        launch from the ground or from a height
  \item explain the \ang{45} rule, use energy, and describe the effect of air resistance
  \item use the uniform velocity and constant acceleration models to predict where and when
        objects will be
  \item plan and judge an investigation that tests the projectile model
\end{goals}

\begin{syllabus}
  \item M5.IQ1.01: analyse projectile motion as a horizontal motion and a vertical motion, with
        two assumptions: constant vertical acceleration and no air resistance
  \item M5.IQ1.02: use the model to derive how the launch velocity, launch angle, maximum
        height, time of flight, final velocity, launch height and range are related
  \item M5.IQ1.03: carry out an investigation that collects primary data to validate these
        relationships
  \item M5.IQ1.04: use the equations of motion (uniform velocity and constant acceleration) to
        solve problems and make predictions
\end{syllabus}

%% ==================================================================
\section{The projectile model}\label{ch1:sec:model}
\covers{M5.IQ1.01, M5.IQ1.02, M5.IQ1.04}

\begin{bigidea}
A thrown ball moves sideways and up-and-down at the same time. The two motions do not
affect each other. Sideways the velocity stays the same, but up and down gravity changes it
at a steady rate.
\end{bigidea}

\begin{keywords}
  \kw{Projectile}{an object that moves through the air after it is thrown or launched, with no
    engine pushing it.}
  \kw{Model}{a simple description of real life that we can calculate with.}
  \kw{Component}{the part of a vector (an arrow quantity) that points in one direction.}
  \kw{Independent}{not affecting each other. (In an experiment, ``independent variable'' means
    something different: see \cref{ch1:sec:experiments}.)}
  \kw{Displacement ($s$)}{the change in position, with a direction. Unit: metre (m).}
\end{keywords}

Kick a football or throw a stone. Once it leaves you, the main force on it is gravity. In
this chapter's model we ignore the push of the air. The ball is a \term{projectile}: an
object that moves through the air after it is thrown or launched, with no engine pushing it.
A rocket with its engine on is not a projectile.

The ball's \term{trajectory} is the path it follows. The path is a \term{parabola}: a curve
that is symmetrical about its highest point. It has this shape for a reason. The sideways
distance grows steadily with time. But the drop caused by gravity grows with the time
squared ($t^2$). A launch from a height uses only part of the curve.

Real motion is complicated, so we use a \term{model}: a simple description of real life
that we can calculate with. The projectile model makes two \term{assumptions} (things we
accept as true to keep the maths simple):
\begin{enumerate}
  \item there is \textbf{zero air resistance} (\term{air resistance} is the push of the air
        on a moving object);
  \item the \textbf{acceleration due to gravity is constant}: $g = \qty{9.8}{m.s^{-2}}$
        straight down, everywhere on the path.
\end{enumerate}

Gravity pulls straight down. So there is no horizontal force (we ignore air resistance).
By Newton's second law, no net force (no overall force) in a direction means no acceleration
in that direction. So the \textbf{horizontal} velocity stays the same for the whole flight.
\textbf{Vertically}, gravity changes the velocity by \qty{9.8}{m.s^{-1}} every second. The two
motions are \term{independent}: they do not affect each other. The only thing they share is
the \emph{time} (\cref{ch1:fig:strobe}).

\begin{sidebyside}[0.47]
  \centering
  % scale 6 cm per metre in both directions; B thrown at 1.5 m/s; positions every 0.1 s
  \begin{tikzpicture}[x=1cm, y=1cm]
    \fill[ground] (-1.9,-5.08) rectangle (4.1,-4.88);
    \draw[hscink, line width=0.6pt] (-1.9,-4.88) -- (4.1,-4.88);
    \foreach \y/\t in {0/0, -0.294/0.1, -1.176/0.2, -2.646/0.3, -4.704/0.4} {
      \draw[guide] (-1.4,\y) -- (3.95,\y);
      \node[diagramlabel, anchor=east, font=\scriptsize, inner sep=1pt] at (-1.45,\y) {\t\,s};
    }
    \foreach \y in {0, -0.294, -1.176, -2.646, -4.704}
      \fill[object] (-0.9,\y) circle (3pt);
    \foreach \x/\y in {0/0, 0.9/-0.294, 1.8/-1.176, 2.7/-2.646, 3.6/-4.704}
      \fill[object, fill=hscseriesB!35] (\x,\y) circle (3pt);
    \draw[velocity] (0.15,0.3) -- (0.9,0.3) node[right, text=hscvel, inner sep=1pt] {$u$};
    \node[diagramlabel, above, inner sep=2pt] at (-0.9,0.12) {\textbf{A}};
    \node[diagramlabel, above, inner sep=2pt] at (0,0.12) {\textbf{B}};
  \end{tikzpicture}
  \captionof{figure}{A is dropped. B is thrown sideways at the same moment.
    Positions every \qty{0.1}{s}.}\label{ch1:fig:strobe}
\nextside
  \begin{itemize}
    \item At every time shown (every \qty{0.1}{s}), A and B are at the \textbf{same height}.
          Moving sideways does not change how B falls.
    \item B moves the \textbf{same distance sideways} in each \qty{0.1}{s}.
    \item Both balls fall further in each \qty{0.1}{s}: they speed up downwards.
    \item So A and B land \textbf{at the same time}.
  \end{itemize}
\end{sidebyside}

To use this idea, we split the launch velocity into two parts. First, a reminder about
vectors.

\begin{refresher}{Vectors and scalars (Year 11)}
A \term{vector} is a quantity with a size and a direction. Velocity, displacement,
acceleration and force are vectors. We draw a vector as an arrow: the length shows the size,
and the arrow shows the direction. A \term{scalar} has a size only. Speed, time, mass and
energy are scalars. \term{Speed} is the size of the velocity. For example,
\qty{5.0}{m.s^{-1}} is a speed. \qty{5.0}{m.s^{-1}} at \ang{53} below the horizontal is a
velocity.
\end{refresher}

We split the launch velocity into a horizontal \term{component} and a vertical component. A
component is the part of a vector that points in one direction.

\begin{toolkit}{Splitting a vector with sin and cos}
\begin{minipage}[t]{0.58\linewidth}\vspace{0pt}\RaggedRight
Draw the velocity $u$ as the longest side of a right-angled triangle. This side is the
\term{hypotenuse}. The angle $\theta$ is measured from the horizontal.
\begin{itemize}
  \item The side next to $\theta$ is $u\cos\theta$, because
        $\cos\theta = \text{adjacent} \div \text{hypotenuse}$.
  \item The side opposite $\theta$ is $u\sin\theta$, because
        $\sin\theta = \text{opposite} \div \text{hypotenuse}$.
  \item To get $u$ from its parts, use Pythagoras: $u = \sqrt{u_x^2 + u_y^2}$.
  \item To get the angle, use inverse tan (the $\tan^{-1}$ key, often SHIFT then tan):
        $\theta = \tan^{-1}(u_y/u_x)$.
\end{itemize}
Use \textbf{degree mode}: $\cos\ang{60}$ must give 0.5.
\end{minipage}\hfill
\begin{minipage}[t]{0.38\linewidth}\vspace{2mm}\centering
  \begin{tikzpicture}[x=1cm, y=1cm]
    \draw[->, line width=0.9pt, hscvel!60] (0,0) -- (2.6,0)
      node[midway, below, font=\small] {$u_x = u\cos\theta$};
    \draw[->, line width=0.9pt, hscvel!60] (2.6,0) -- (2.6,1.5)
      node[midway, right, font=\small] {$u_y = u\sin\theta$};
    \draw[velocity] (0,0) -- (2.6,1.5) node[midway, above left, text=hscvel, inner sep=1pt] {$u$};
    \draw[hscink, line width=0.5pt] (2.4,0) -- (2.4,0.2) -- (2.6,0.2);
    \draw[hscink, line width=0.5pt] (0.75,0) arc (0:30:0.75);
    \node[font=\small] at (14:1.0) {$\theta$};
  \end{tikzpicture}
\end{minipage}
\end{toolkit}

\begin{formulacard}[notonsheet]{Components of the launch velocity}
  \cardformula{u_x = u\cos\theta \qquad u_y = u\sin\theta}
  \begin{symbols}
    \sym{$u$}{launch speed}{\si{m.s^{-1}}}
    \sym{$\theta$}{launch angle, measured up from the horizontal (work in degrees)}{\si{\degree} (SI unit: rad)}
    \sym{$u_x$, $u_y$}{horizontal and vertical parts of the launch velocity}{\si{m.s^{-1}}}
  \end{symbols}
  \inwords{the launch velocity splits into a sideways part and an upwards part.}
  \whentouse{at the start of almost every projectile question with an angle.}
  \whennottouse{the angle is from the vertical: then $u_x = u\sin\theta$, $u_y = u\cos\theta$.}
  \rearranged{$u = \sqrt{u_x^2 + u_y^2}$ \qquad $\theta = \tan^{-1}\!\left(\dfrac{u_y}{u_x}\right)$}
  \tensecond{\qty{20}{m.s^{-1}} at \ang{30}: $u_x = \qty{17.3}{m.s^{-1}}$,
    $u_y = \qty{10.0}{m.s^{-1}}$.}
  \derive{in the triangle $\cos\theta = \dfrac{u_x}{u}$ and $\sin\theta = \dfrac{u_y}{u}$.
    Multiply both sides of each by $u$.}
  \trap{radian mode gives $20\cos 30 = 3.09$, not 17.3. Each component is smaller than $u$.}
\end{formulacard}

To follow the motion in each direction, we use the equations of motion from Year 11.

\begin{refresher}{The equations of motion (Year 11)}
For straight-line motion with constant acceleration, three equations link $u$ (initial
velocity), $v$ (final velocity), $a$ (acceleration), $t$ (time) and the \term{displacement}
$s$ (the change in position, with a direction). For a projectile, use them once for each
direction:\par
\smallskip
\begin{tabular}{@{}L{32mm}L{40mm}L{56mm}@{}}
  \toprule
  & \colhead{Horizontal ($x$)} & \colhead{Vertical ($y$)}\\
  \midrule
  Acceleration & $a_x = 0$ & $a_y = -g = \qty{-9.8}{m.s^{-2}}$\\
  Velocity at time $t$ & $v_x = u_x$ (never changes) & $v_y = u_y + a_y t$\\
  Displacement at time $t$ & $s_x = u_x t$ & $s_y = u_y t + \frac{1}{2}a_y t^2$\\
  Equation with no $t$ & (not needed) & $v_y^2 = u_y^2 + 2a_y s_y$\\
  \bottomrule
\end{tabular}
\end{refresher}

The \term{origin} is the point we measure every position from ($x = 0$, $y = 0$). In this book
the origin is the launch point, and \textbf{up and forwards are positive}. The size of
gravity's acceleration is $g = \qty{9.8}{m.s^{-2}}$, a positive number. It points down, so
$a_y = -g = \qty{-9.8}{m.s^{-2}}$ for the whole flight, including the top.
\Cref{ch1:fig:graphs} shows the motion as graphs.

\begin{toolkit}{Gradient of a straight-line graph}
Gradient $= \dfrac{\text{rise}}{\text{run}} = \dfrac{\Delta y}{\Delta x}$. Use two points far
apart on the line. For data, use the \term{line of best fit}: the straight line that passes as
close as possible to all the points. The gradient's unit is the $y$ unit divided by the $x$
unit. Example: in the $v_y$--$t$ graph of \cref{ch1:fig:graphs}, $v_y$ changes from
\qty{12.0}{m.s^{-1}} at $t = 0$ to \qty{-12.0}{m.s^{-1}} at $t = \qty{2.449}{s}$. So the
gradient is $\dfrac{-12.0 - 12.0}{2.449 - 0} = \qty{-9.80}{m.s^{-2}}$: the acceleration $a_y$.
\end{toolkit}

\begin{figure}
  \centering
  \includegraphics{ch1_motiongraphs.pdf}
  \caption{Graphs against time for a ball launched on level ground with
    $u_x = \qty{9.0}{m.s^{-1}}$ and $u_y = \qty{12.0}{m.s^{-1}}$. Top row: horizontal motion.
    Bottom row: vertical motion.}\label{ch1:fig:graphs}
\end{figure}

Here are the three equations you will use in every question. To use them, you often need to
rearrange them, and at the end you round the answer.

\begin{toolkit}{Rearranging a formula, and rounding the answer}
\textbf{Rearranging.} To find an unknown, get it alone on one side of the equals sign.
Whatever you do to one side, do the same to the other side. Undo each operation with its
opposite: $+$ with $-$, $\times$ with $\div$, a square with a square root. Example: find $t$
from $v = u + at$. (1) Subtract $u$ from both sides: $v - u = at$. (2) Divide both sides by
$a$: $t = \dfrac{v - u}{a}$. Remember: a negative divided by a negative is positive, so
$(-1.25) \div (-4.9) = +0.255$.

\smallskip
\textbf{Rounding.} \term{Significant figures} are the digits that carry meaning: \num{1.25}
has 3 and \num{4.2} has 2. Keep all the calculator digits until the end. Then round to the
number of significant figures of the least precise value in the question. Do not count
$g = 9.8$: it is a constant from the data sheet. HSC markers usually accept 2 or 3 significant
figures.
\end{toolkit}

\begin{formulacard}{The equations of motion}
  \cardformula{v = u + at \qquad s = ut + \frac{1}{2}at^2 \qquad v^2 = u^2 + 2as}
  \begin{symbols}
    \sym{$u$, $v$}{initial velocity, final velocity}{\si{m.s^{-1}}}
    \sym{$a$}{acceleration (must be constant)}{\si{m.s^{-2}}}
    \sym{$s$, $t$}{displacement (it has a sign); time}{\si{m}; \si{s}}
  \end{symbols}
  \inwords{each equation links four of the five quantities $u$, $v$, $a$, $s$, $t$.}
  \whentouse{vertically with $a_y = \qty{-9.8}{m.s^{-2}}$; horizontally with $a_x = 0$,
    which gives $s_x = u_x t$.}
  \whennottouse{do not mix $x$ values and $y$ values in one equation. Do not use these
    equations if the acceleration changes (for example, with large air resistance).}
  \rearranged{$t = \dfrac{v - u}{a}$ \qquad $s = \dfrac{v^2 - u^2}{2a}$ \qquad from rest:
    $t = \sqrt{\dfrac{2s}{a}}$}
  \tensecond{$u_y = \qty{12}{m.s^{-1}}$ up; after \qty{0.50}{s},
    $v_y = 12 + (-9.8)(0.50) = \qty{7.1}{m.s^{-1}}$ (still rising).}
  \trap{with up positive, $a_y$ is always negative and $s_y$ is negative below the launch
    point. Choose the sign of a square root from the direction of motion.}
\end{formulacard}

\subsection{The simplest case: a horizontal launch}
A marble rolls off a table, or a skateboarder rolls off a ledge. At the edge the velocity is
horizontal, so $u_y = 0$. The fall is then like a drop from rest: ball B in
\cref{ch1:fig:strobe} falls exactly like ball A. So the time in the air depends only on the
height, not on the speed.

% book assembler: this card must not start at the foot of a page. (The consistency edit
% replaced the assembler's \clearpage by a conditional \Needspace: after the syllabus box was
% added, the \clearpage left a nearly empty page.)
\Needspace{14\baselineskip}
\begin{formulacard}[notonsheet]{Horizontal launch from a height}
  \cardformula{t = \sqrt{\frac{2h}{g}} \qquad R = u\sqrt{\frac{2h}{g}}}
  \begin{symbols}
    \sym{$t$}{time of flight}{\si{s}}
    \sym{$h$}{launch height above the landing point (positive)}{\si{m}}% layout QA: one line
    \sym{$g$}{size of gravity's acceleration, \num{9.8}}{\si{m.s^{-2}}}
    \sym{$u$}{horizontal launch speed}{\si{m.s^{-1}}}
    \sym{$R$}{range, from the point under the launch point}{\si{m}}
  \end{symbols}
  \inwords{the fall time depends only on the height, not on the speed.}
  \whentouse{the launch velocity is horizontal ($u_y = 0$), for example off a table, a ledge
    or a cliff.}
  \whennottouse{the launch is at an angle ($u_y \neq 0$): solve the quadratic in
    \cref{ch1:sec:height} instead.}
  \rearranged{$h = \tfrac{1}{2}g t^2$ \qquad $u = \dfrac{R}{t}$}
  \tensecond{\qty{3.0}{m.s^{-1}} off a \qty{0.80}{m} table:
    $t = \sqrt{2 \times 0.80/9.8} = \qty{0.404}{s}$, $R = 3.0 \times 0.404 = \qty{1.21}{m}$.}
  \derive{start from the sheet equation for the vertical motion.
    % layout QA: slightly tighter lines so the card's last line (Trap) stays on its page
    \setlength{\jot}{1pt}%
    \begin{align*}
      -h &= 0 \times t + \tfrac{1}{2}(-g)t^2 && \text{$s = ut + \tfrac12at^2$ with $s_y = -h$,
        $u_y = 0$, $a_y = -g$}\\
      h &= \tfrac{1}{2}g t^2 && \text{multiply both sides by $-1$}\\
      t^2 &= \frac{2h}{g} && \text{multiply both sides by 2, divide both sides by $g$}\\
      t &= \sqrt{\frac{2h}{g}} && \text{square root (time is positive)}\\
      R &= u_x t = u\sqrt{\frac{2h}{g}} && \text{horizontal: $s_x = u_x t$}
    \end{align*}}
  \trap{use $u_y = 0$, not $u_y = u$. Here $h$ is a positive height: do not put in $-h$,
    because the minus sign is already used in the derivation.}
\end{formulacard}

\textbf{How the launch height changes the flight.} A higher launch point gives more time to
fall, so the time of flight and the range both grow. Four times the height gives two times
the time, because $\sqrt{4} = 2$. For example, a fall of \qty{1.25}{m} takes \qty{0.505}{s},
and a fall of \qty{5.0}{m} takes \qty{1.01}{s}. At the same launch speed the range also
doubles. The landing speed grows too: $v = \sqrt{u^2 + 2gh}$ (\cref{ch1:sec:energy}).

\begin{workedexample}{A skateboarder rolls off a ledge}\label{ch1:we:ledge}
A skateboarder rolls off a horizontal ledge \qty{1.25}{m} high at \qty{4.2}{m.s^{-1}}.
How long is she in the air, and how far from the ledge does she land?
\given{$u_x = \qty{4.2}{m.s^{-1}}$, $u_y = 0$ (horizontal launch), $s_y = \qty{-1.25}{m}$,
  $a_y = \qty{-9.8}{m.s^{-2}}$}
\find{the time $t$ (vertical motion), then $s_x$ (horizontal motion)}
\formula{$s_y = u_y t + \frac{1}{2}a_y t^2$ \quad and \quad $s_x = u_x t$}
\substitute{\eqlines{\qty{-1.25}{m} &= 0 + \tfrac{1}{2}(\qty{-9.8}{m.s^{-2}})\,t^2 \quad\text{(vertical; $u_y = 0$, so $u_y t = 0$)}\\
  \qty{-1.25}{m} &= (\qty{-4.9}{m.s^{-2}})\,t^2 \quad\text{($\tfrac12 \times (-9.8) = -4.9$)}\\
  t^2 &= \frac{\qty{-1.25}{m}}{\qty{-4.9}{m.s^{-2}}} = \qty{0.2551}{s^2} \quad\text{(divide both sides by $-4.9$)}\\
  t &= \sqrt{\qty{0.2551}{s^2}} = \qty{0.5051}{s} \quad\text{(square root; time is positive)}\\
  s_x &= \qty{4.2}{m.s^{-1}} \times \qty{0.5051}{s} = \qty{2.121}{m} \quad\text{(horizontal; the same $t$)}}}
\answer{\final{$t = \qty{0.505}{s}$} (3 significant figures, like $h = \qty{1.25}{m}$) and
  \final{$s_x = \qty{2.1}{m}$} (2 significant figures, like $u = \qty{4.2}{m.s^{-1}}$)}
\checkstep{put $t$ back in: $\tfrac{1}{2} \times 9.8 \times 0.5051^2 = \qty{1.25}{m}$, the
  ledge height \cmark. In about half a second at \qty{4.2}{m.s^{-1}} she moves about
  \qty{2}{m} \cmark. The horizontal-launch card gives the same $t$ \cmark.}
\end{workedexample}

\begin{workedexample}{Where is the ball after 1.5 s?}\label{ch1:we:position}
A ball is thrown from the ground at \qty{16}{m.s^{-1}} at \ang{40} above the horizontal.
Find its position and its vertical velocity \qty{1.5}{s} later. Is it rising or falling?
\given{$u = \qty{16}{m.s^{-1}}$, $\theta = \ang{40}$, $t = \qty{1.5}{s}$, $a_y = \qty{-9.8}{m.s^{-2}}$}
\find{$s_x$, $s_y$ and $v_y$}
\formula{$u_x = u\cos\theta$, $u_y = u\sin\theta$; \ $s_x = u_x t$; \
  $s_y = u_y t + \frac{1}{2}a_y t^2$; \ $v_y = u_y + a_y t$}
\substitute{\eqlines{u_x &= 16\cos\ang{40} = \qty{12.26}{m.s^{-1}} \quad\text{(split $u$ into parts; degree mode)}\\
  u_y &= 16\sin\ang{40} = \qty{10.285}{m.s^{-1}} \quad\text{(keep the extra digits)}}\par
  \eqlines{s_x &= \qty{12.26}{m.s^{-1}} \times \qty{1.5}{s} = \qty{18.39}{m} \quad\text{(horizontal: $s_x = u_x t$)}\\
  s_y &= (\qty{10.285}{m.s^{-1}})(\qty{1.5}{s}) + \tfrac{1}{2}(\qty{-9.8}{m.s^{-2}})(\qty{2.25}{s^2})\\
    &\qquad\text{(vertical; $1.5^2 = 2.25$)}\\
    &= \qty{15.43}{m} - \qty{11.03}{m} = \qty{4.40}{m} \quad\text{($\tfrac12 \times (-9.8) \times 2.25 = -11.03$)}\\
  v_y &= \qty{10.285}{m.s^{-1}} + (\qty{-9.8}{m.s^{-2}})(\qty{1.5}{s}) = \qty{-4.42}{m.s^{-1}}\\
    &\qquad\text{(vertical velocity: $v = u + at$)}}}
\answer{\final{\qty{18}{m} forwards, \qty{4.4}{m} up}; \final{$v_y = \qty{-4.4}{m.s^{-1}}$}:
  the minus sign means it is \textbf{falling}.}
\checkstep{the top comes at $t = 10.285/9.8 = \qty{1.05}{s}$, before \qty{1.5}{s}, so
  falling makes sense \cmark. $s_y > 0$: still in the air \cmark.}
\end{workedexample}

\begin{mistakes}
\begin{itemize}
  \item \xmark\ ``The throw gives a forward force.'' \cmark\ After launch the only force in the
        model is gravity (air resistance is ignored).
  \item \xmark\ Using sin to find the horizontal part. \cmark\ When the angle is measured from
        the horizontal, $u_x = u\cos\theta$.
  \item \xmark\ Using $u_y = u$ for a horizontal launch. \cmark\ For a horizontal launch, $u_y = 0$.
  \item \xmark\ Thinking a faster horizontal launch stays in the air longer. \cmark\ The fall
        time depends only on the height.
  \item \xmark\ Using $a_y = +\qty{9.8}{m.s^{-2}}$, or a positive $s_y$ for a fall. \cmark\ With
        up positive, $a_y = \qty{-9.8}{m.s^{-2}}$, and $s_y$ is negative below the launch point.
\end{itemize}
\end{mistakes}

\begin{examtip}
\begin{itemize}
  \item ``Explain why the horizontal velocity stays constant'' (2 marks).
        (1) The only force is gravity (air resistance is ignored), and gravity acts vertically.
        So there is no horizontal force. (2) By Newton's second law
        ($\vec{F}_\mathrm{net} = m\vec{a}$), no horizontal force means no horizontal
        acceleration. So $v_x$ stays constant.
  \item Graph questions (\cref{ch1:fig:graphs}): $s_x$--$t$ is a straight line through the
        origin; $s_y$--$t$ is an upside-down parabola; $v_x$--$t$ is flat; $v_y$--$t$ is a
        straight line with gradient \qty{-9.8}{m.s^{-2}}; $a_y$--$t$ is flat at
        \qty{-9.8}{m.s^{-2}}.
\end{itemize}
\end{examtip}

\begin{quickcheck}
  \item Find $u_x$ and $u_y$ for (a) \qty{15}{m.s^{-1}} at \ang{60} above the horizontal;
        (b) a dart thrown \emph{upwards} at \qty{8.0}{m.s^{-1}} at \ang{20} \emph{to the vertical}.
  \item A stone is thrown horizontally at \qty{6.0}{m.s^{-1}} from a cliff \qty{20}{m}
        above the water. Find its time in the air and how far from the cliff it lands.
  \item Ball A is dropped. At the same moment ball B is thrown horizontally from the same
        height. Which lands first? Explain.
  \item A marble rolls horizontally off a table. The table is replaced by one 4 times as high,
        and the marble leaves it at the same speed. What happens to (a) its time in the air and
        (b) its range? Explain.
\end{quickcheck}

%% ==================================================================
\section{Height, time of flight, range and final velocity}\label{ch1:sec:height}
\covers{M5.IQ1.02, M5.IQ1.04}

\begin{bigidea}
At the top of the path the vertical velocity is zero for an instant. This gives the
maximum height and, on level ground, the time of flight. The range is the horizontal
velocity multiplied by the time of flight.
\end{bigidea}

\begin{keywords}
  \kw{Maximum height ($h_\mathrm{max}$)}{the greatest height above the launch point.}
  \kw{Time of flight ($t_\mathrm{flight}$)}{the time from launch until landing.}
  \kw{Range ($R$)}{the horizontal distance from the launch point to the landing point.}
  \kw{Launch height ($h$)}{how far the launch point is above the landing point.}
  \kw{Final velocity}{the velocity just before landing: a speed and a direction.}
\end{keywords}

\subsection{The top of the flight}
In \cref{ch1:fig:path}, the horizontal arrow $v_x$ has the same length everywhere. The
vertical arrow $v_y$ shrinks on the way up, is zero at the top, then grows downwards. At the
top the ball still moves at $v_x$, and its acceleration is still \qty{9.8}{m.s^{-2}} down.
On level ground the path is \term{symmetrical}: the way down mirrors the way up.

\begin{figure}
  \centering
  \begin{tikzpicture}[x=1cm, y=1cm,
      comp/.style={->, line width=0.8pt, color=hscvel!55}]
    \draw[hscink, line width=0.6pt] (-0.4,0) -- (9.4,0);
    \draw[hscmuted!60, line width=1.1pt, domain=0:2.449, samples=80, variable=\t]
      plot ({0.37*9*\t},{0.37*(12*\t-4.9*\t*\t)});
    \draw[comp] (0,0) -- (1.08,0);
    \draw[comp] (0,0) -- (0,1.44) node[left, font=\footnotesize, text=hscvel] {$u_y$};
    \draw[velocity] (0,0) -- (1.08,1.44) node[above, text=hscvel, inner sep=1pt] {$u$};
    \node[font=\footnotesize, text=hscvel, below] at (0.6,0) {$u_x$};
    \draw[comp] (2.039,2.039) -- (3.119,2.039);
    \draw[comp] (2.039,2.039) -- (2.039,2.759);
    \draw[velocity] (2.039,2.039) -- (3.119,2.759);
    \draw[velocity] (4.078,2.718) -- (5.158,2.718)
      node[right, text=hscvel, inner sep=2pt, font=\footnotesize] {$v = v_x$, \ $v_y = 0$};
    \draw[accel] (4.078,2.56) -- (4.078,1.80)
      node[below, text=hscacc, font=\footnotesize, inner sep=1pt] {$a = g$ (down)};
    \draw[comp] (6.117,2.039) -- (7.197,2.039);
    \draw[comp] (6.117,2.039) -- (6.117,1.319);
    \draw[velocity] (6.117,2.039) -- (7.197,1.319);
    \draw[comp] (8.156,0) -- (9.236,0);
    \draw[comp] (8.156,0) -- (8.156,-1.44) node[left, font=\footnotesize, text=hscvel] {$v_y$};
    \draw[velocity] (8.156,0) -- (9.236,-1.44) node[right, text=hscvel, inner sep=1pt] {$v$};
    \foreach \p in {(0,0),(2.039,2.039),(4.078,2.718),(6.117,2.039),(8.156,0)}
      \fill[object] \p circle (3.2pt);
    \draw[guide] (7.5,2.718) -- (10.05,2.718);
    \draw[displacement, <->] (9.9,0) -- (9.9,2.718)
      node[midway, right, text=hscink, font=\small] {$h_\mathrm{max}$};
    \draw[<->, line width=0.6pt, hscink] (0,-1.7) -- (8.156,-1.7)
      node[midway, fill=white, font=\small, inner sep=2pt] {range $R$};
    \node[font=\footnotesize, text=hscmuted] at (2.0,-0.5) {time up $= t_\mathrm{top}$};
    \node[font=\footnotesize, text=hscmuted] at (6.1,-0.5) {time down $= t_\mathrm{top}$};
  \end{tikzpicture}
  \caption{Velocity along a flight on level ground. Light arrows: $v_x$ and $v_y$. Dark
    arrows: the full velocity, along the path. The acceleration is $g$ downwards at every
    point; it is drawn only at the top.}\label{ch1:fig:path}
\end{figure}

These three results come straight from the equations of motion.

\begin{formulacard}[notonsheet]{The top of the flight, and the time of flight}
  \cardformula{t_\mathrm{top} = \frac{u_y}{g} \qquad h_\mathrm{max} = \frac{u_y^2}{2g}
    \qquad t_\mathrm{flight} = \frac{2u_y}{g}}
  \begin{symbols}
    \sym{$t_\mathrm{top}$, $t_\mathrm{flight}$}{time to the top; time from launch to landing}{\si{s}}
    \sym{$h_\mathrm{max}$}{greatest height above the launch point}{\si{m}}
    \sym{$u_y$}{vertical part of the launch velocity, $u\sin\theta$}{\si{m.s^{-1}}}
    \sym{$g$}{size of gravity's acceleration, \num{9.8}}{\si{m.s^{-2}}}
  \end{symbols}
  \inwords{all three depend only on $u_y$. Double $u_y$: the times double, and the height
    becomes 4 times bigger.}
  \whentouse{$t_\mathrm{top}$ and $h_\mathrm{max}$: any launch that goes upwards
    ($u_y > 0$), from any height. $t_\mathrm{flight}$: launch and landing at the same height.}
  \whennottouse{$t_\mathrm{flight} = 2u_y/g$ when the landing point is lower than the launch
    point: solve the quadratic below instead. Do not use $h_\mathrm{max}$ as the height above
    the ground after a launch from a height: add the launch height $h$ first.}
  \rearranged{$u_y = g\,t_\mathrm{top} = \sqrt{2g\,h_\mathrm{max}} = \dfrac{g\,t_\mathrm{flight}}{2}$.
    With $u_y = u\sin\theta$: \ $h_\mathrm{max} = \dfrac{u^2\sin^2\theta}{2g}$ \ and \
    $t_\mathrm{flight} = \dfrac{2u\sin\theta}{g}$. ($\sin^2\theta$ means $(\sin\theta)^2$.)
    For example, \qty{20}{m.s^{-1}} at \ang{30} gives $h_\mathrm{max} = \qty{5.10}{m}$ and
    $t_\mathrm{flight} = \qty{2.04}{s}$.}
  \tensecond{$u_y = \qty{14}{m.s^{-1}}$: $t_\mathrm{top} = \qty{1.43}{s}$,
    $h_\mathrm{max} = \dfrac{14^2}{2 \times 9.8} = \qty{10.0}{m}$, $t_\mathrm{flight} = \qty{2.86}{s}$.}
  \derive{put $a_y = -g$ into the equations on the formula sheet. At the top $v_y = 0$; at
    landing $s_y = 0$.
    \begin{align*}
      0 &= u_y - g\,t_\mathrm{top} && \text{$v = u + at$ with $v_y = 0$, $a_y = -g$}\\
      g\,t_\mathrm{top} &= u_y && \text{add $g\,t_\mathrm{top}$ to both sides}\\
      t_\mathrm{top} &= \frac{u_y}{g} && \text{divide both sides by $g$}\\[4pt]
      0 &= u_y^2 - 2g\,h_\mathrm{max} && \text{$v^2 = u^2 + 2as$ with $v_y = 0$, $s_y = h_\mathrm{max}$}\\
      2g\,h_\mathrm{max} &= u_y^2 && \text{add $2g\,h_\mathrm{max}$ to both sides}\\
      h_\mathrm{max} &= \frac{u_y^2}{2g} && \text{divide both sides by $2g$}\\[4pt]
      0 &= u_y t - \tfrac{1}{2}g t^2 && \text{$s = ut + \tfrac12 at^2$ with $s_y = 0$}\\
      0 &= t\left(u_y - \tfrac{1}{2}g t\right) && \text{take out the common factor $t$}\\
      0 &= u_y - \tfrac{1}{2}g t && \text{$t = 0$ is the launch, so the bracket is zero}\\
      \tfrac{1}{2}g t &= u_y && \text{add $\tfrac12 gt$ to both sides}\\
      t_\mathrm{flight} &= \frac{2u_y}{g} = 2t_\mathrm{top} && \text{multiply both sides by $\tfrac{2}{g}$}
    \end{align*}}
  \trap{use $u_y$, not $u$. In these three formulas $g = 9.8$ (a positive number), because the
    minus sign of the acceleration is already built in. Do not put in $-9.8$. Also,
    $t_\mathrm{top}$ is only half of the flight.}
\end{formulacard}

\subsection{Range and the \texorpdfstring{\ang{45}}{45 degree} rule}
Point a garden hose a little upwards and the water lands further away. Tilt it more and it
goes further still. But point it almost straight up and the water lands close to you. So
there is a best angle in between. A low launch moves forwards fast, but it lands quickly. A
steep launch stays up for a long time, but it moves forwards slowly. The range needs both:
horizontal speed and time in the air. \Cref{ch1:fig:angles} shows launches at one speed and
five angles.

\begin{figure}
  \centering
  \includegraphics{ch1_angles.pdf}
  \caption{Launches at \qty{20}{m.s^{-1}} (level ground, no air resistance). Lines of the
    same style land at the same point.}\label{ch1:fig:angles}
\end{figure}

Speeds in questions are sometimes given in \si{km.h^{-1}}. Every formula in this book needs
\si{m.s^{-1}}.

\begin{toolkit}{Changing units}
$\qty{1}{km} = \qty{1000}{m}$ and $\qty{1}{h} = \qty{3600}{s}$. So
$\qty{1}{km.h^{-1}} = \dfrac{\qty{1000}{m}}{\qty{3600}{s}} = \dfrac{1}{3.6}\,\si{m.s^{-1}}$.
\begin{itemize}
  \item \si{km.h^{-1}} to \si{m.s^{-1}}: divide by 3.6. Example: $54 \div 3.6 = \qty{15}{m.s^{-1}}$.
  \item \si{mm} (millimetre) to \si{m}, and \si{ms} (millisecond) to \si{s}: divide by 1000.
        Example: $\qty{16.0}{mm} = \qty{0.0160}{m}$.
\end{itemize}
\end{toolkit}

\begin{toolkit}{The double-angle rule}
For any angle, $2\sin\theta\cos\theta = \sin 2\theta$. Test with \ang{30}:
$2 \times 0.5 \times 0.866 = 0.866 = \sin\ang{60}$ \cmark. The largest value of $\sin$ is 1, at
\ang{90}. So $\sin 2\theta$ is largest when $2\theta = \ang{90}$, that is, when
$\theta = \ang{45}$. On a calculator, type $\sin(2 \times 40)$, not $2\sin 40$.
\end{toolkit}

\begin{formulacard}[notonsheet]{Range on level ground}
  \cardformula{R = \frac{u^2\sin 2\theta}{g}}
  \begin{symbols}
    \sym{$R$}{range: horizontal distance to the landing point}{\si{m}}
    \sym{$u$}{launch speed}{\si{m.s^{-1}}}
    \sym{$\theta$}{launch angle above the horizontal}{\si{\degree} (SI unit: rad)}
  \end{symbols}
  \inwords{double the speed and the range becomes 4 times bigger; for one speed, the range is
    largest at \ang{45}.}
  \whentouse{launch and landing at the same height, no air resistance.}
  \whennottouse{launches from a height: find $t$ from $s_y$ first, then $R = u_x t$.}
  \rearranged{$u = \sqrt{\dfrac{Rg}{\sin 2\theta}}$ \qquad $R_\mathrm{max} = \dfrac{u^2}{g}$
    at \ang{45} \qquad $\theta = \tfrac{1}{2}\sin^{-1}\!\left(\dfrac{Rg}{u^2}\right)$, and also
    $\ang{90} - \theta$ (two angles give the same range). Example: \qty{18}{m.s^{-1}} and
    $R = \qty{30}{m}$ give \ang{32.6} or \ang{57.4}.}
  \tensecond{\qty{10}{m.s^{-1}} at \ang{45}: $R = \dfrac{10^2 \times \sin\ang{90}}{9.8} = \qty{10.2}{m}$.}
  \derive{range $=$ horizontal velocity $\times$ time of flight:
    \begin{align*}
      R &= u_x\,t_\mathrm{flight} = u\cos\theta \times \frac{2u\sin\theta}{g}
        && \text{$s_x = u_x t$ at landing}\\
        &= \frac{2u^2\sin\theta\cos\theta}{g} && \text{multiply: $u \times u = u^2$}\\
        &= \frac{u^2\,(2\sin\theta\cos\theta)}{g} = \frac{u^2\sin 2\theta}{g}
        && \text{the double-angle rule}
    \end{align*}}
  \trap{level ground only. Change \si{km.h^{-1}} to \si{m.s^{-1}} first (divide by 3.6).}
\end{formulacard}

\begin{itemize}
  \item \textbf{\ang{45} gives the longest range}, because $\sin 2\theta = 1$ when
        $2\theta = \ang{90}$. This is true \emph{only} on level ground with no air resistance.
  \item \textbf{\term{Complementary angles}} (two angles that add to \ang{90}) give the same
        range. Example: \ang{30} and \ang{60}. $\sin(2 \times \ang{30}) = \sin\ang{60} = 0.866$
        and $\sin(2 \times \ang{60}) = \sin\ang{120} = 0.866$. The steeper one goes higher and
        stays up longer.
  \item \textbf{Double the speed} at the same angle (level ground): the range and the height
        become 4 times bigger; the time of flight becomes 2 times bigger.
\end{itemize}

\subsection{Launching from a height}
A ball thrown from a cliff lands below its launch point, so the flight is not symmetrical.
Set up the question as in \cref{ch1:fig:cliff}. Then follow these steps:
\begin{enumerate}
  \item Put the origin at the launch point, with up positive. Then $s_y = -h$ at landing.
  \item Find $u_x$ and $u_y$.
  \item Solve $-h = u_y t + \frac{1}{2}a_y t^2$ for $t$. This is a \term{quadratic equation}:
        an equation with a $t^2$ term, a $t$ term and a number.
  \item Find the range: $R = u_x t$.
\end{enumerate}
For a horizontal launch, $u_y = 0$ and the equation is simpler: use the card ``Horizontal
launch from a height'' in \cref{ch1:sec:model}.

A higher launch point gives the projectile more time in the air, so a longer range. The
maximum height \emph{above the launch point} does not change, because
$h_\mathrm{max} = u_y^2/2g$ depends only on $u_y$.

If the target is \emph{above} the launch point (for example, a basketball hoop), $s_y$ is
positive. The quadratic then has two positive roots. The smaller one is on the way up; the
larger one is on the way down. Choose the one the question describes.

\begin{figure}
  \centering
  \begin{tikzpicture}[x=1cm, y=1cm,
      comp/.style={->, line width=0.8pt, color=hscvel!55}]
    \fill[ground] (-1.4,-3.9) rectangle (0,0);
    \fill[ground] (0,-4.1) rectangle (6.6,-3.9);
    \draw[hscink, line width=0.6pt] (-1.4,0) -- (0,0) -- (0,-3.9) -- (6.6,-3.9);
    \draw[axisline] (0,0) -- (0,1.05) node[above, diagramlabel] {$+y$ (up)};
    \draw[axisline] (0,0) -- (1.45,0) node[right, diagramlabel] {$+x$};
    \node[diagramlabel, above left, inner sep=2pt] at (0,0) {O};
    \draw[hscmuted!60, line width=1.1pt, domain=0:2.3327, samples=70, variable=\t]
      plot ({0.26*8.660*\t},{0.26*(5*\t-4.9*\t*\t)});
    \draw[velocity] (0,0) -- (1.126,0.65) node[above, text=hscvel, inner sep=1pt] {$u$};
    \draw[hscink, line width=0.5pt] (0.62,0) arc (0:30:0.62);
    \node[font=\footnotesize] at (14:0.85) {$\theta$};
    \draw[guide] (2.1,0) -- (6.3,0);
    \draw[displacement, <->] (6.15,0) -- (6.15,-3.9);
    \node[diagramlabel, align=left, anchor=west] at (6.25,-1.95) {$h$\\(so $s_y = -h$)};
    \fill[object] (5.253,-3.9) circle (3.2pt);
    \draw[<->, line width=0.6pt, hscink] (0,-4.45) -- (5.253,-4.45)
      node[midway, below, font=\small] {range $R$ (from the foot of the cliff)};
    \node[diagramlabel, anchor=west] at (8.0,-0.35) {final velocity:};
    \draw[comp] (8.3,-0.9) -- (9.426,-0.9) node[midway, above, font=\footnotesize, text=hscvel] {$v_x$};
    \draw[comp] (9.426,-0.9) -- (9.426,-3.222) node[midway, right, font=\footnotesize, text=hscvel] {$v_y$};
    \draw[velocity] (8.3,-0.9) -- (9.426,-3.222) node[midway, left, text=hscvel] {$v$};
    \draw[hscink, line width=0.5pt] (8.8,-0.9) arc (0:-64.1:0.5);
    \node[font=\footnotesize] at (8.95,-1.35) {$\alpha$};
  \end{tikzpicture}
  \caption{A launch from a height. The origin O is the launch point, so the landing point
    has $s_y = -h$. Right: the final velocity is made from $v_x$ and $v_y$.}
    \label{ch1:fig:cliff}
\end{figure}

\begin{toolkit}{Solving a quadratic equation}
Write the equation as $At^2 + Bt + C = 0$. (We use capital letters, because $a$ already
means acceleration.) Then
\[ t = \frac{-B \pm \sqrt{B^2 - 4AC}}{2A}. \]
Example: $4.9t^2 - 3.0t - 2.0 = 0$ has $A = 4.9$, $B = -3.0$, $C = -2.0$. First,
$B^2 - 4AC = (-3.0)^2 - 4(4.9)(-2.0) = 9.0 + 39.2 = 48.2$, and $-B = +3.0$. So
\[ t = \frac{3.0 \pm \sqrt{48.2}}{9.8} = \frac{3.0 \pm 6.943}{9.8}
   = \qty{1.01}{s} \ \text{or}\ {\qty{-0.402}{s}}. \]
A \term{root} is a solution of the equation. Keep the \textbf{positive} root. The negative
one is a time before the launch.
\end{toolkit}

\subsection{Final velocity}
The \term{final velocity} is the velocity just before landing. Think of catching a ball thrown
to you from far away. It arrives moving forwards and downwards at the same time. Its final
velocity combines both parts, like the arrow on the right of \cref{ch1:fig:cliff}. The
horizontal part is still $v_x = u_x$. The vertical part comes from $v_y = u_y + a_y t$ or
$v_y^2 = u_y^2 + 2a_y s_y$.

\begin{formulacard}[notonsheet]{Speed and direction from components}
  \cardformula{v = \sqrt{v_x^2 + v_y^2} \qquad \tan\alpha = \frac{\lvert v_y\rvert}{v_x}}
  \begin{symbols}
    \sym{$v$}{speed (size of the velocity)}{\si{m.s^{-1}}}
    \sym{$v_x$, $v_y$}{horizontal ($= u_x$) and vertical velocity at that moment}{\si{m.s^{-1}}}
    \sym{$\lvert v_y\rvert$}{the size of $v_y$ without its sign, e.g. $\lvert -4.0\rvert = 4.0$}{\si{m.s^{-1}}}
    \sym{$\alpha$}{angle of the velocity below (or above) the horizontal}{\si{\degree} (SI unit: rad)}
  \end{symbols}
  \inwords{Pythagoras gives the speed; inverse tan gives the direction.}
  \whentouse{at landing, or at any moment, once you know $v_x$ and $v_y$.}
  \whennottouse{do not put $u_y$ in place of $v_y$. The vertical velocity changes during the
    flight, so find $v_y$ at that moment first.}
  \rearranged{$\alpha = \tan^{-1}\!\left(\dfrac{\lvert v_y\rvert}{v_x}\right)$ \qquad $v_x = v\cos\alpha$}
  \tensecond{$v_x = \qty{3.0}{m.s^{-1}}$, $v_y = \qty{-4.0}{m.s^{-1}}$: $v = \qty{5.0}{m.s^{-1}}$
    at $\tan^{-1}(4.0/3.0) = \ang{53}$ below the horizontal.}
  \derive{$v_x$ and $v_y$ are the short sides of a right-angled triangle; $v$ is the long side
    (\cref{ch1:fig:cliff}). Use Pythagoras, then $\tan = \text{opposite} \div \text{adjacent}$.}
  \trap{a velocity needs a direction: write ``\ang{53} below the horizontal''.}
\end{formulacard}

On \textbf{level ground} the final speed equals the launch speed. The final angle equals the
launch angle, but below the horizontal. Why: at landing $s_y = 0$, so
$v_y^2 = u_y^2 + 2a_y \times 0 = u_y^2$. So $v_y = -u_y$: the same size, but moving down. And
$v_x = u_x$ never changes. So the velocity triangle at landing is the launch triangle flipped
upside down.

\textbf{From a height}, the best angle for range is \emph{less} than \ang{45}. Even a low throw
gets extra time in the air from the drop. So it pays to put more of the speed into the
horizontal part, $u_x$. For the shot put in \cref{ch1:we:shotput} (below), the best angle is
about \ang{42}.

\begin{workedexample}{A soccer kick}\label{ch1:we:soccer}
A soccer ball is kicked from level ground at \qty{21}{m.s^{-1}} at \ang{32} above the
horizontal. Find its maximum height, its time of flight and its range.
\given{$u = \qty{21}{m.s^{-1}}$, $\theta = \ang{32}$, level ground, $g = \qty{9.8}{m.s^{-2}}$}
\find{$h_\mathrm{max}$, $t_\mathrm{flight}$ and $R$}
\formula{$h_\mathrm{max} = \dfrac{u_y^2}{2g}$; \ $t_\mathrm{flight} = \dfrac{2u_y}{g}$; \
  $R = u_x\,t_\mathrm{flight}$}
\substitute{\eqlines{u_x &= 21\cos\ang{32} = \qty{17.81}{m.s^{-1}} \quad\text{(split $u$ into parts)}\\
  u_y &= 21\sin\ang{32} = \qty{11.13}{m.s^{-1}}}\par
  \eqlines{h_\mathrm{max} &= \frac{(\qty{11.13}{m.s^{-1}})^2}{2 \times \qty{9.8}{m.s^{-2}}} \quad\text{(top card)}\\
    &= \frac{\qty{123.9}{m^2.s^{-2}}}{\qty{19.6}{m.s^{-2}}} = \qty{6.32}{m}
    \quad\text{(square $u_y$, then divide by $2g$)}\\
  t_\mathrm{flight} &= \frac{2 \times \qty{11.13}{m.s^{-1}}}{\qty{9.8}{m.s^{-2}}} = \qty{2.271}{s}
    \quad\text{(level ground)}\\
  R &= \qty{17.81}{m.s^{-1}} \times \qty{2.271}{s} = \qty{40.45}{m}
    \quad\text{($R = u_x \times t_\mathrm{flight}$)}}}
\answer{\final{$h_\mathrm{max} = \qty{6.3}{m}$}, \final{$t_\mathrm{flight} = \qty{2.3}{s}$},
  \final{$R = \qty{40}{m}$} (2 significant figures, like the data: 40.45 rounds to 40, and
  here the 0 counts)}
\checkstep{the range card gives $\dfrac{21^2 \sin\ang{64}}{9.8} = \qty{40.4}{m}$ \cmark.}
\end{workedexample}

\begin{workedexample}{Does the ball clear the wall?}\label{ch1:we:wall}
A ball is kicked from the ground at \qty{79.2}{km.h^{-1}} at \ang{30} above the horizontal. A
wall \qty{4.0}{m} high is \qty{30.0}{m} away. Does the ball go over the wall?
\given{$u = \qty{79.2}{km.h^{-1}} = 79.2 \div 3.6 = \qty{22.0}{m.s^{-1}}$, $\theta = \ang{30}$,
  $s_x = \qty{30.0}{m}$ at the wall, $a_y = \qty{-9.8}{m.s^{-2}}$}
\find{the height $s_y$ when $s_x = \qty{30.0}{m}$; compare it with \qty{4.0}{m}}
\formula{$s_x = u_x t$; \ $s_y = u_y t + \frac{1}{2}a_y t^2$}
\substitute{\eqlines{u_x &= 22.0\cos\ang{30} = \qty{19.05}{m.s^{-1}} \quad\text{(split $u$ into parts)}\\
  u_y &= 22.0\sin\ang{30} = \qty{11.00}{m.s^{-1}}\\
  t &= \frac{s_x}{u_x} = \frac{\qty{30.0}{m}}{\qty{19.05}{m.s^{-1}}} = \qty{1.5746}{s}
    \quad\text{(time to reach the wall)}}\par
  \eqlines{s_y &= (\qty{11.00}{m.s^{-1}})(\qty{1.5746}{s}) + \tfrac{1}{2}(\qty{-9.8}{m.s^{-2}})(\qty{1.5746}{s})^2\\
    &\qquad\text{(vertical: the height at that time)}\\
    &= \qty{17.32}{m} - \qty{12.15}{m} = \qty{5.17}{m} \quad\text{($4.9 \times 1.5746^2 = 12.15$)}}}
\answer{\final{$s_y = \qty{5.2}{m}$}, which is more than \qty{4.0}{m}: \textbf{yes}, the ball
  goes over the wall, by about \qty{1.2}{m}.}
\checkstep{$t_\mathrm{flight} = 2 \times 11.00/9.8 = \qty{2.24}{s}$. This is more than
  \qty{1.57}{s}, so the ball is still in the air at the wall \cmark.}
\end{workedexample}

\begin{workedexample}{A shot put released above the ground}\label{ch1:we:shotput}
A shot put leaves the hand \qty{1.90}{m} above the ground at \qty{12.0}{m.s^{-1}} at
\ang{41.0} above the horizontal. Find (a) the time of flight, (b) the range and (c) the
velocity just before it lands.
\given{$u = \qty{12.0}{m.s^{-1}}$, $\theta = \ang{41.0}$, $s_y = \qty{-1.90}{m}$ (origin at
  the hand, up positive), $a_y = \qty{-9.8}{m.s^{-2}}$}
\find{$t$; \ $R = u_x t$; \ the speed $v$ and the angle $\alpha$}
\formula{$s_y = u_y t + \frac{1}{2}a_y t^2$; \ $v_y^2 = u_y^2 + 2a_y s_y$; \
  $v = \sqrt{v_x^2 + v_y^2}$; \ $\tan\alpha = \dfrac{\lvert v_y\rvert}{v_x}$}
\substitute{\eqlines{u_x &= 12.0\cos\ang{41.0} = \qty{9.057}{m.s^{-1}} \quad\text{(split $u$ into parts)}\\
  u_y &= 12.0\sin\ang{41.0} = \qty{7.873}{m.s^{-1}}}\par
  \textbf{(a)} Vertical motion, $s_y = u_yt + \tfrac12 a_yt^2$:\par
  \eqlines{\qty{-1.90}{m} &= (\qty{7.873}{m.s^{-1}})\,t + \tfrac{1}{2}(\qty{-9.8}{m.s^{-2}})\,t^2\\
  -1.90 &= 7.873\,t - 4.9\,t^2 \quad\text{($\tfrac12 \times (-9.8) = -4.9$)}\\
  0 &= 7.873\,t - 4.9\,t^2 + 1.90 \quad\text{(add 1.90 to both sides)}\\
  0 &= 4.9\,t^2 - 7.873\,t - 1.90 \quad\text{(multiply every term by $-1$)}}\par
  Quadratic formula with $A = 4.9$, $B = -7.873$, $C = -1.90$:\par
  \eqlines{B^2 - 4AC &= (-7.873)^2 - 4(4.9)(-1.90)\\
    &= 61.98 + 37.24 = 99.22 \quad\text{(minus $\times$ minus $=$ plus)}\\
  t &= \frac{-B \pm \sqrt{99.22}}{2A} = \frac{7.873 \pm 9.961}{9.8} \quad\text{($-B = +7.873$)}\\
  t &= \frac{17.834}{9.8} = \qty{1.820}{s} \quad\text{or}\quad \frac{-2.088}{9.8} = \qty{-0.213}{s}}\par
  Keep the positive root: $t = \qty{1.820}{s}$.\par
  \textbf{(b)} Horizontal motion:
  $R = u_x t = \qty{9.057}{m.s^{-1}} \times \qty{1.820}{s} = \qty{16.48}{m}$.\par
  \textbf{(c)} Vertical velocity from $v^2 = u^2 + 2as$, then Pythagoras and inverse tan:\par
  \eqlines{v_y^2 &= (\qty{7.873}{m.s^{-1}})^2 + 2(\qty{-9.8}{m.s^{-2}})(\qty{-1.90}{m})\\
    &= \qty{61.98}{m^2.s^{-2}} + \qty{37.24}{m^2.s^{-2}} = \qty{99.22}{m^2.s^{-2}}\\
  v_y &= -\sqrt{\qty{99.22}{m^2.s^{-2}}} = \qty{-9.961}{m.s^{-1}} \quad\text{(negative: it moves down)}\\
  v &= \sqrt{(\qty{9.057}{m.s^{-1}})^2 + (\qty{9.961}{m.s^{-1}})^2} = \qty{13.46}{m.s^{-1}}\\
  \alpha &= \tan^{-1}(9.961 \div 9.057) = \ang{47.7} \quad\text{(opposite $\div$ adjacent)}}}
\answer{(a) \final{\qty{1.82}{s}} \quad (b) \final{\qty{16.5}{m}}\\
  (c) \final{\qty{13.5}{m.s^{-1}} at \ang{47.7} below the horizontal} \ (3 significant
  figures, like the data 12.0, 41.0 and 1.90)}
\checkstep{the level-ground card would give only \qty{1.61}{s}; the shot falls an extra
  \qty{1.90}{m}, so it stays up longer and lands faster than \qty{12.0}{m.s^{-1}} \cmark.}
\end{workedexample}

\begin{mistakes}
\begin{itemize}
  \item \xmark\ ``At the top, the velocity and the acceleration are zero.'' \cmark\ Only
        $v_y = 0$. The acceleration is still \qty{9.8}{m.s^{-2}} down.
  \item \xmark\ Using $t_\mathrm{flight} = 2u_y/g$ or $R = u^2\sin 2\theta/g$ for a launch from a
        height. \cmark\ Solve the quadratic with $s_y = -h$, and keep the positive root.
  \item \xmark\ ``A higher angle always goes further.'' \cmark\ On level ground, the range gets
        smaller for angles above \ang{45}.
  \item \xmark\ Using the maximum height to decide whether a ball clears a wall. \cmark\ Find
        the time at the wall from $s_x = u_x t$, then $s_y$ at that time.
  \item \xmark\ Giving only the speed for a ``final velocity''. \cmark\ Also give the direction,
        e.g. ``\ang{48} below the horizontal''.
\end{itemize}
\end{mistakes}

\begin{examtip}
\begin{itemize}
  \item ``Show that the range on level ground is greatest at \ang{45}'' (3 marks).
        (1) Get $t_\mathrm{flight}$ from $s_y = 0$. (2) Use $R = u_x t_\mathrm{flight}$ and the
        double-angle rule to get $R = u^2\sin 2\theta/g$. (3) $\sin 2\theta$ is at most 1, at
        $2\theta = \ang{90}$. State the conditions: level ground, no air resistance.
  \item Final velocity (3 marks): give $v_x$, then $v_y$ with its sign, then $v$ \emph{and}
        its direction.
  \item ``Does it clear the wall?'': find $s_y$ at the wall's distance. Finish with a clear
        yes or no.
\end{itemize}
\end{examtip}

\begin{quickcheck}
  \item A ball is launched from level ground at \qty{25}{m.s^{-1}} at \ang{50}. Find its
        maximum height, its time of flight and its speed at the top.
  \item A ball is kicked from level ground at \qty{64.8}{km.h^{-1}}. Find its range at
        \ang{45}, \ang{30} and \ang{60}.
  \item Water leaves a hose nozzle at \ang{40} above the horizontal. It lands \qty{6.5}{m}
        away, at the same height as the nozzle. Find the speed of the water leaving the hose.
  \item A stone is thrown at \qty{12}{m.s^{-1}} at \ang{25} above the horizontal from the top
        of an \qty{18}{m} cliff. Find its time in the air, how far from the foot of the cliff
        it lands, and its velocity just before it hits the sea.
  \item A rugby ball is kicked from the ground at \qty{20}{m.s^{-1}} at \ang{35}. The
        crossbar is \qty{3.0}{m} high and \qty{30}{m} away. Does the ball go over it?
\end{quickcheck}

%% ==================================================================
\section{Energy and air resistance}\label{ch1:sec:energy}
\covers{M5.IQ1.01, M5.IQ1.02}

\begin{bigidea}
With no air resistance, kinetic plus potential energy stays the same during the flight.
This gives the speed at any height without the time or the angle. Real air takes energy
away, so real projectiles go lower and shorter than the model predicts.
\end{bigidea}

\begin{keywords}
  \kw{Kinetic energy ($K$)}{energy an object has because it moves.}
  \kw{Gravitational potential energy ($U$)}{energy stored because of height.}
  \kw{Conservation of energy}{the total energy stays the same.}
  \kw{Drag}{another word for air resistance.}
\end{keywords}

\begin{refresher}{Energy (Year 11)}
\term{Kinetic energy} is the energy of motion: $K = \frac{1}{2}mv^2$. \term{Gravitational
potential energy} is the energy of height. Rising by $\Delta h$ adds $\Delta U = mg\Delta h$;
falling by $\Delta h$ takes the same amount away. When gravity is the only force that acts
(no air resistance), $K + U$ stays the same. This is \term{conservation of energy}.
\end{refresher}

\begin{sidebyside}[0.44]
  \centering
  % launch at 60 deg on level ground: K_top / K_launch = cos^2(60) = 0.25
  \begin{tikzpicture}[x=1cm, y=1cm]
    \draw[hscink, line width=0.6pt] (-0.3,0) -- (4.5,0);
    \fill[hscseriesA!75] (0,0) rectangle (0.8,2.4);
    \fill[hscseriesA!75] (1.8,0) rectangle (2.6,0.6);
    \fill[hscseriesB!70] (1.8,0.6) rectangle (2.6,2.4);
    \fill[hscseriesA!75] (3.6,0) rectangle (4.4,2.4);
    \node[font=\footnotesize, text=white] at (0.4,1.2) {$K$};
    \node[font=\footnotesize, text=white] at (2.2,0.3) {$K$};
    \node[font=\footnotesize, text=white] at (2.2,1.5) {$U$};
    \node[font=\footnotesize, text=white] at (4.0,1.2) {$K$};
    \node[font=\scriptsize, text=hscink] at (0.4,2.62) {\qty{100}{\percent}};
    \node[font=\scriptsize, text=hscink, anchor=east] at (1.75,0.3) {\qty{25}{\percent}};
    \node[font=\scriptsize, text=hscink, anchor=east] at (1.75,1.5) {\qty{75}{\percent}};
    \node[font=\scriptsize, text=hscink] at (4.0,2.62) {\qty{100}{\percent}};
    \node[font=\footnotesize, text=hscink, below] at (0.4,-0.05) {launch};
    \node[font=\footnotesize, text=hscink, below] at (2.2,-0.05) {top};
    \node[font=\footnotesize, text=hscink, below] at (4.0,-0.05) {landing};
  \end{tikzpicture}
  \captionof{figure}{Energy of a ball launched at \ang{60} on level ground, as a percentage of
    the total. Blue: $K$. Orange: $U$.}\label{ch1:fig:energybars}
\nextside
  On the way up, kinetic energy turns into potential energy. On the way down, it turns back
  into kinetic energy.\par\smallskip
  At the top of an angled flight, $K$ is smallest but \emph{not} zero. The ball still moves
  sideways at $v_x$, so $K_\mathrm{top} = \frac{1}{2}mu_x^2$. In \cref{ch1:fig:energybars},
  $K$ at the top is \qty{25}{\percent} of the total, because $\cos^2\ang{60} = 0.25$.
\end{sidebyside}

Energy has no direction, so it gives only the \emph{speed}. For the angle, use $v_x = u_x$
and $\cos\alpha = v_x/v$ (adjacent $\div$ hypotenuse in the velocity triangle of
\cref{ch1:fig:cliff}).

\begin{formulacard}[combine={$K = \frac{1}{2}mv^2$ and $\Delta U = mg\Delta h$}]{Energy conservation for a projectile}
  \cardformula{\frac{1}{2}mv_2^2 - \frac{1}{2}mv_1^2 = mg(h_1 - h_2)}
  \begin{symbols}
    \sym{$m$}{mass of the projectile}{\si{kg}}
    \sym{$v_1$, $v_2$}{speed at point 1 and at point 2}{\si{m.s^{-1}}}
    \sym{$h_1$, $h_2$}{height at point 1 and at point 2, above any level you choose as zero.
      Only $h_1 - h_2$ (the $\Delta h$ in $\Delta U = mg\Delta h$) matters}{\si{m}}
  \end{symbols}
  \inwords{the kinetic energy gained equals the potential energy lost.}
  \whentouse{you need the speed at a known height. You do not need the time or the angle to
    find the speed.}
  \whennottouse{you need the direction, or air resistance matters.}
  \rearranged{$v_2 = \sqrt{v_1^2 + 2g(h_1 - h_2)}$ \qquad landing $h$ below the launch:
    $v = \sqrt{u^2 + 2gh}$}
  \tensecond{thrown at \qty{6.0}{m.s^{-1}}, lands \qty{5.0}{m} lower:
    $v = \sqrt{6.0^2 + 2 \times 9.8 \times 5.0} = \qty{11.6}{m.s^{-1}}$.}
  \derive{the kinetic energy gained is $\frac{1}{2}mv_2^2 - \frac{1}{2}mv_1^2$ (from
    $K = \frac12mv^2$). The potential energy lost is $mg(h_1 - h_2)$ (from
    $\Delta U = mg\Delta h$). With no air resistance these are equal.
    \begin{align*}
      v_2^2 - v_1^2 &= 2g(h_1 - h_2) && \text{divide every term by $\tfrac12 m$: the mass cancels}\\
      v_2^2 &= v_1^2 + 2g(h_1 - h_2) && \text{add $v_1^2$ to both sides}
    \end{align*}}
  \trap{$h_1 - h_2$ is positive when the ball ends lower, and negative when it ends higher.
    Energy gives only the speed: for the direction you still need $v_x = u_x$ and
    $\cos\alpha = v_x/v$.}
\end{formulacard}

\begin{workedexample}{A stone from a bridge}\label{ch1:we:bridge}
A stone is thrown from a bridge \qty{32}{m} above a river at \qty{14}{m.s^{-1}}, in an
unknown direction. Find its speed as it reaches the water. Ignore air resistance.
\given{$v_1 = \qty{14}{m.s^{-1}}$; it falls $h_1 - h_2 = \qty{32}{m}$; $g = \qty{9.8}{m.s^{-2}}$}
\find{the speed $v_2$ at the water}
\formula{$\frac{1}{2}mv_2^2 - \frac{1}{2}mv_1^2 = mg(h_1 - h_2)$, so
  $v_2 = \sqrt{v_1^2 + 2g(h_1 - h_2)}$}
\substitute{\eqlines{v_2 &= \sqrt{(\qty{14}{m.s^{-1}})^2 + 2 \times \qty{9.8}{m.s^{-2}} \times \qty{32}{m}}
    \quad\text{(mass cancelled)}\\
  &= \sqrt{\qty{196}{m^2.s^{-2}} + \qty{627.2}{m^2.s^{-2}}} \quad\text{($14^2 = 196$; \ $2 \times 9.8 \times 32 = 627.2$)}\\
  &= \sqrt{\qty{823.2}{m^2.s^{-2}}} = \qty{28.69}{m.s^{-1}} \quad\text{(add, then square root)}}}
\answer{\final{$v_2 = \qty{29}{m.s^{-1}}$} (2 significant figures), whatever the angle.}
\checkstep{dropped from rest it would reach $\sqrt{627.2} = \qty{25}{m.s^{-1}}$; thrown, it
  should be faster \cmark.}
\end{workedexample}

\begin{workedexample}{Speed and direction at a height}\label{ch1:we:energyangle}
A ball is launched from the ground at \qty{15}{m.s^{-1}} at \ang{60} above the horizontal.
Find its speed and the angle of its velocity when it is \qty{6.0}{m} above the ground.
\given{$u = v_1 = \qty{15}{m.s^{-1}}$, $\theta = \ang{60}$, $h_1 = 0$, $h_2 = \qty{6.0}{m}$}
\find{the speed $v_2$ and the angle $\alpha$ to the horizontal}
\formula{$v_2^2 = v_1^2 + 2g(h_1 - h_2)$; \ $v_x = u\cos\theta$; \ $\cos\alpha = \dfrac{v_x}{v_2}$}
\substitute{\eqlines{v_2^2 &= (\qty{15}{m.s^{-1}})^2 + 2 \times \qty{9.8}{m.s^{-2}} \times (0 - \qty{6.0}{m})\\
    &\qquad\text{($h_1 - h_2 = 0 - 6.0 = \qty{-6.0}{m}$, because it rises)}\\
  &= \qty{225}{m^2.s^{-2}} - \qty{117.6}{m^2.s^{-2}} = \qty{107.4}{m^2.s^{-2}}\\
    &\qquad\text{($15^2 = 225$; \ $2 \times 9.8 \times 6.0 = 117.6$)}\\
  v_2 &= \sqrt{\qty{107.4}{m^2.s^{-2}}} = \qty{10.36}{m.s^{-1}} \quad\text{(square root)}\\
  v_x &= \qty{15}{m.s^{-1}} \times \cos\ang{60} = \qty{7.50}{m.s^{-1}} \quad\text{($v_x = u_x$ never changes)}\\
  \alpha &= \cos^{-1}(7.50 \div 10.36) = \ang{43.6} \quad\text{(adjacent $\div$ hypotenuse)}}}
\answer{\final{\qty{10}{m.s^{-1}} at \ang{44} to the horizontal} (2 significant figures, like
  the data; in 10 the 0 counts). It passes \qty{6.0}{m} twice: \ang{44} \emph{above} the
  horizontal going up, \ang{44} \emph{below} it coming down.}
\checkstep{it does reach \qty{6.0}{m}: $h_\mathrm{max} = (15\sin\ang{60})^2/19.6 = \qty{8.6}{m}$
  \cmark. Energy alone cannot tell up from down.}
\end{workedexample}

\subsection{Air resistance: where the model fails}
In real life \term{drag} (air resistance) acts on every projectile. It points opposite to the
velocity and grows quickly with speed. \Cref{ch1:fig:drag} compares the two paths.
\begin{itemize}
  \item The horizontal velocity gets smaller during the flight, so the \textbf{range is shorter}.
  \item On the way up drag adds to gravity, so the \textbf{maximum height is lower}.
  \item The path is \textbf{not symmetrical}: the way down is steeper and takes longer.
  \item The ball \textbf{lands slower and steeper}, and the best angle is \textbf{below \ang{45}}.
\end{itemize}

\begin{figure}
  \centering
  \includegraphics{ch1_drag.pdf}
  \caption{\qty{30}{m.s^{-1}} at \ang{45}. With drag (computer model) the range is reduced from
    about \qty{92}{m} to \qty{69}{m}, and the maximum height from \qty{23}{m} to
    \qty{19}{m}.}\label{ch1:fig:drag}
\end{figure}

The model works well for small, dense objects at low speeds over short distances. (\term{Dense}
means heavy for its size.) Examples are a steel ball, a shot put and a stone. At the same
speed, a steel ball and a paper ball of the same size feel about the same drag. But the steel
ball is much heavier, so drag hardly changes its motion.

The model works badly for light or large objects, such as a feather or a ball of paper. It
also works badly for fast objects, such as a golf ball after a drive.

Over very long distances, gravity gets weaker with height and the Earth curves away. This
leads to orbits (\cref{ch4:chap}). In Module 6, a charged particle between charged plates
follows the same kind of parabola.

\begin{mistakes}
\begin{itemize}
  \item \xmark\ ``Kinetic energy is zero at the top.'' \cmark\ Only for a throw straight up.
        For an angled throw, $K_\mathrm{top} = \frac{1}{2}mu_x^2$, which is not zero.
  \item \xmark\ Using energy to find a direction. \cmark\ Energy gives only the speed. Find the
        angle with $\cos\alpha = v_x/v$.
  \item \xmark\ ``Air resistance makes gravity weaker.'' \cmark\ Gravity is unchanged; drag is
        an extra force.
\end{itemize}
\end{mistakes}

\begin{examtip}
``Explain how air resistance affects the trajectory'' or ``Evaluate the projectile model''
(3--5 marks). Full marks need:
\begin{itemize}
  \item the force and its direction: drag, opposite to the velocity, bigger at higher speed;
  \item a link from drag to each effect: shorter range, lower height, steeper and longer
        fall (a sketch of both paths helps);
  \item for ``evaluate'', a judgement: good for dense, slow objects over short distances;
        poor for light or fast ones.
\end{itemize}
\end{examtip}

\begin{quickcheck}
  \item A ball is thrown at \qty{9.0}{m.s^{-1}}, in an unknown direction, from a balcony
        \qty{1.5}{m} above the ground. Find its speed as it hits the ground.
  \item A ball is launched at \ang{40} above the horizontal. What fraction of its launch
        kinetic energy does it still have at the top? Explain.
  \item Describe three ways air resistance changes the trajectory of a tennis ball.
\end{quickcheck}

%% ==================================================================
% Section added by the syllabus coverage audit (M5.IQ1.04 was thin: the uniform-velocity
% model was never named, and there was no one-dimensional or two-object prediction).
\section{Predicting motion: uniform velocity and constant acceleration}\label{ch1:sec:predict}
\covers{M5.IQ1.04}

\begin{bigidea}
Two simple models describe motion in a straight line: uniform velocity and constant
acceleration. A projectile uses both at once. Write one equation for each object and each
direction, with the same time $t$, and you can predict where and when things will be.
\end{bigidea}

\begin{keywords}
  \kw{Uniform velocity}{a velocity that does not change: the same speed in the same direction.
    The acceleration is zero.}
  \kw{Constant acceleration}{an acceleration that does not change, so the velocity changes by
    the same amount every second.}
  \kw{Free fall}{motion when gravity is the only force acting.}
  \kw{Prediction}{a result you calculate from a model before you measure it.}
\end{keywords}

A model is useful when it makes \term{predictions}: results we calculate before we measure
them. Then we can test the model. This section shows how to choose a model for each part of a
motion. It also shows how to join the parts to predict where and when objects meet.

\textbf{Model 1: uniform velocity.} An object has \term{uniform velocity} when its velocity
does not change: the same speed in the same direction. This happens when the net force on it is
zero (Newton's first law). The object moves equal distances in equal times. Examples: a car on
cruise control on a straight, flat road; a puck sliding on smooth ice; the \emph{horizontal}
motion of a projectile. Its displacement--time graph is a straight line. The gradient of this
line is the velocity.

\textbf{Model 2: constant acceleration.} When the net force stays the same, the acceleration
stays the same ($\vec{F}_\mathrm{net} = m\vec{a}$). This is \term{constant acceleration}: the
velocity changes by the same amount every second. Examples: a stone in \term{free fall}
(gravity is the only force on it), a car braking steadily, and the \emph{vertical} motion of a
projectile. Use the three equations of motion (\cref{ch1:sec:model}). The velocity--time graph
is a straight line, and its gradient is the acceleration.

\textbf{Which model?} Ask one question for each direction: is there a net force this way?
\begin{itemize}
  \item No net force: use \textbf{uniform velocity}.
  \item A net force that stays the same: use \textbf{constant acceleration}.
  \item A force that changes (for example, large air resistance): neither model is exact.
\end{itemize}

\begin{formulacard}[combine={$s = ut + \frac{1}{2}at^2$ and $v = u + at$, with $a = 0$}]{Uniform velocity}
  \cardformula{s = ut \qquad v = u}
  \begin{symbols}
    \sym{$s$}{displacement in the direction of motion}{\si{m}}
    \sym{$u$}{the velocity, which does not change}{\si{m.s^{-1}}}
    \sym{$t$}{time}{\si{s}}
    \sym{$v$}{velocity at time $t$ (the same as $u$)}{\si{m.s^{-1}}}
  \end{symbols}
  \inwords{with no acceleration, the displacement grows by the same amount every second.}
  \whentouse{there is no net force in that direction: the horizontal motion of a projectile
    ($s_x = u_x t$), or a vehicle moving at a steady speed in a straight line.}
  \whennottouse{the speed or the direction changes, for example the vertical motion of a
    projectile or a car that brakes. Use the three constant-acceleration equations instead.}
  \rearranged{$u = \dfrac{s}{t}$ \qquad $t = \dfrac{s}{u}$ \qquad starting at position $x_0$:
    \ $x = x_0 + ut$}
  \tensecond{a cart moving at \qty{0.60}{m.s^{-1}} for \qty{0.40}{s} moves
    $0.60 \times 0.40 = \qty{0.24}{m}$.}
  \derive{put $a = 0$ into two equations on the formula sheet.
    \begin{align*}
      s &= ut + \tfrac{1}{2}(0)t^2 = ut && \text{$s = ut + \tfrac12 at^2$ with $a = 0$}\\
      v &= u + (0)t = u && \text{$v = u + at$ with $a = 0$}
    \end{align*}}
  \trap{do not use $s = ut$ for a falling object, because there $a \neq 0$. Change
    \si{km.h^{-1}} to \si{m.s^{-1}} first.}
\end{formulacard}

\textbf{Throwing straight up.} A ball thrown straight up is a projectile with $u_x = 0$. Only
the vertical model is needed: constant acceleration, with $a_y = \qty{-9.8}{m.s^{-2}}$. The
model makes three predictions that you can test.
\begin{itemize}
  \item The time up equals the time down.
  \item The ball comes back to the launch point at its launch speed, but moving down.
  \item Below the top, the ball passes each height twice: once going up and once coming down.
        So the equation for the time has two positive answers.
\end{itemize}

Graphs are models too. The next toolkit shows how to read a displacement from a
velocity--time graph.

\begin{toolkit}{Area under a velocity--time graph}
On a velocity--time graph, the gradient is the acceleration. The \textbf{area} between the line
and the time axis is the \textbf{displacement}. Area above the axis is positive, and area below
the axis is negative. For uniform velocity the area is a rectangle: $s = ut$. For a velocity
that falls steadily to zero, the area is a triangle: $\tfrac12 \times \text{base} \times
\text{height}$. Example: in the $v_y$--$t$ graph of \cref{ch1:fig:graphs}, $v_y$ falls from
\qty{12.0}{m.s^{-1}} to zero in \qty{1.2245}{s}. The area of this triangle is
$\tfrac12 \times \qty{1.2245}{s} \times \qty{12.0}{m.s^{-1}} = \qty{7.35}{m}$. This is the
maximum height. (Check: $h_\mathrm{max} = 12.0^2 \div 19.6 = \qty{7.35}{m}$ \cmark.) On the way
down, the triangle below the axis has the same area, so at landing the vertical displacement
is zero.
\end{toolkit}

\textbf{Objects launched from moving vehicles.} An object let go from a moving vehicle keeps the
vehicle's velocity at that moment. This is \term{inertia}: the tendency of an object to keep its
velocity. In \cref{ch1:fig:cart}, a cart moves at a uniform velocity of \qty{2.0}{m.s^{-1}}. A
ball is tossed straight up from it (as seen from the cart) at \qty{4.9}{m.s^{-1}}.
\begin{itemize}
  \item Seen from the cart, the ball goes straight up and straight down.
  \item Seen from the ground, the ball also moves forwards at \qty{2.0}{m.s^{-1}}. Uniform
        velocity sideways plus constant acceleration downwards gives a parabola.
  \item The ball lands back in the cart after \qty{1.0}{s}. By then both have moved
        \qty{2.0}{m}. (This needs no air resistance, and a cart that keeps a uniform velocity.)
\end{itemize}
In the same way, a parcel dropped from a plane that flies at a uniform velocity stays directly
below the plane until it lands.

\begin{figure}
  \centering
  % 2.4 cm per metre in both directions. Ball: x = 2.0t, y = 4.9t - 4.9t^2 (relative to the
  % top of the cart), positions every 0.25 s (numbers.py). Faint carts: t = 0.25, 0.50, 0.75 s.
  \begin{tikzpicture}[x=2.4cm, y=2.4cm]
    \fill[ground] (-0.45,-0.44) rectangle (2.55,-0.38);
    \draw[hscink, line width=0.6pt] (-0.45,-0.38) -- (2.55,-0.38);
    \foreach \x/\op in {0.5/0.25, 1.0/0.25, 1.5/0.25, 0/1, 2.0/1} {
      \draw[object, opacity=\op] (\x-0.2,-0.3) rectangle (\x+0.2,-0.02);
      \fill[hscink, opacity=\op] (\x-0.12,-0.32) circle (0.055);
      \fill[hscink, opacity=\op] (\x+0.12,-0.32) circle (0.055);
    }
    \draw[hscmuted!60, line width=1pt, domain=0:1, samples=60, variable=\t]
      plot ({2.0*\t},{4.9*\t-4.9*\t*\t});
    \foreach \x/\y in {0.5/0.9188, 1.0/1.225, 1.5/0.9188} \draw[guide] (\x,\y) -- (\x,-0.02);
    \foreach \x/\y/\lab in {0/0/0, 0.5/0.9188/0.25, 1.0/1.225/0.50, 1.5/0.9188/0.75, 2.0/0/1.00} {
      \fill[object, fill=hscseriesB!60] (\x,\y) circle (3pt);
      \node[diagramlabel, font=\scriptsize, anchor=south west, inner sep=2pt] at (\x,\y) {\lab\,s};
    }
    \draw[velocity] (-0.08,0.02) -- (-0.08,0.5) node[left, text=hscvel, font=\footnotesize,
      align=right, inner sep=5pt] {\qty{4.9}{m.s^{-1}} up\\(seen from the cart)};
    \draw[velocity] (-0.2,-0.6) -- (0.4,-0.6) node[right, text=hscvel, font=\footnotesize,
      inner sep=2pt] {cart: \qty{2.0}{m.s^{-1}}, uniform velocity};
    \draw[<->, line width=0.5pt, hscink] (0,-0.82) -- (2.0,-0.82)
      node[midway, below, font=\footnotesize] {\qty{2.0}{m} in \qty{1.0}{s}};
  \end{tikzpicture}
  \caption{A ball tossed straight up from a cart that moves at a uniform velocity. Positions
    every \qty{0.25}{s}; the faint carts show where the cart is at those times. The ball is
    always directly above the cart, and it lands back in it.}\label{ch1:fig:cart}
\end{figure}

\Needspace{6\baselineskip}
\textbf{How to make a prediction.} Follow these six steps.
\begin{enumerate}
  \item Draw a sketch. Choose an origin and a positive direction, and mark them.
  \item For each object and each direction, choose the model: uniform velocity or constant
        acceleration.
  \item Write a position equation for each one. Use one clock: the same $t$, counted from the
        same moment.
  \item Use the condition in the question (for example, ``they are at the same place'') to get
        one equation in $t$. Solve it.
  \item Put $t$ back in to find the place. Check that $t$ is positive and that the object is
        still in the air.
  \item State the assumptions: no air resistance, and a constant $g$.
\end{enumerate}

\begin{workedexample}{A juggler's ball}\label{ch1:we:juggler}
A juggler throws a ball straight up from her hand at \qty{7.0}{m.s^{-1}}. (a) How long is it
until the ball is back at her hand? (b) At what times is the ball \qty{2.0}{m} above her hand?
(c) What is its velocity at each of those times?
\given{$u_y = \qty{7.0}{m.s^{-1}}$ (up), $a_y = \qty{-9.8}{m.s^{-2}}$; origin at the hand, up
  positive}
\find{(a) $t$ when $s_y = 0$ \quad (b) $t$ when $s_y = \qty{2.0}{m}$ \quad (c) $v_y$ at those
  times}
\formula{(a) $t_\mathrm{flight} = \dfrac{2u_y}{g}$ \quad (b) $s_y = u_y t + \frac{1}{2}a_y t^2$
  \quad (c) $v_y = u_y + a_y t$}
\substitute{\textbf{(a)} Launch and landing at the same height (the hand):\par
  \eqlines{t &= \frac{2 \times \qty{7.0}{m.s^{-1}}}{\qty{9.8}{m.s^{-2}}} = \qty{1.429}{s}}\par
  \textbf{(b)} Vertical motion, constant acceleration:\par
  \eqlines{\qty{2.0}{m} &= (\qty{7.0}{m.s^{-1}})\,t + \tfrac{1}{2}(\qty{-9.8}{m.s^{-2}})\,t^2\\
  2.0 &= 7.0\,t - 4.9\,t^2 \quad\text{($\tfrac12 \times (-9.8) = -4.9$)}\\
  0 &= 7.0\,t - 4.9\,t^2 - 2.0 \quad\text{(subtract 2.0 from both sides)}\\
  0 &= 4.9\,t^2 - 7.0\,t + 2.0 \quad\text{(multiply every term by $-1$)}}\par
  Quadratic formula with $A = 4.9$, $B = -7.0$, $C = 2.0$:\par
  \eqlines{B^2 - 4AC &= (-7.0)^2 - 4(4.9)(2.0) = 49.0 - 39.2 = 9.80\\
  t &= \frac{-B \pm \sqrt{9.80}}{2A} = \frac{7.0 \pm 3.130}{9.8} \quad\text{($-B = +7.0$)}\\
  t &= \frac{3.870}{9.8} = \qty{0.3949}{s} \quad\text{or}\quad \frac{10.130}{9.8} = \qty{1.0337}{s}}\par
  Both roots are positive: the first is on the way up, the second on the way down.\par
  \textbf{(c)} Velocity from $v = u + at$ at each time:\par
  \eqlines{v_y &= \qty{7.0}{m.s^{-1}} + (\qty{-9.8}{m.s^{-2}})(\qty{0.3949}{s}) = +\qty{3.13}{m.s^{-1}}\\
  v_y &= \qty{7.0}{m.s^{-1}} + (\qty{-9.8}{m.s^{-2}})(\qty{1.0337}{s}) = \qty{-3.13}{m.s^{-1}}}}
\answer{(a) \final{\qty{1.4}{s}} \quad (b) \final{\qty{0.39}{s} and \qty{1.0}{s}} \quad
  (c) \final{\qty{3.1}{m.s^{-1}} up}, then \final{\qty{3.1}{m.s^{-1}} down}
  (2 significant figures, like the data)}
\checkstep{the top is at $h_\mathrm{max} = 7.0^2 \div 19.6 = \qty{2.5}{m}$. This is above
  \qty{2.0}{m}, so two times make sense \cmark. The two times add up to
  $0.3949 + 1.0337 = \qty{1.429}{s}$, the whole flight: the motion is symmetrical \cmark.
  Energy gives the same speed: $\sqrt{7.0^2 - 2 \times 9.8 \times 2.0} = \qty{3.13}{m.s^{-1}}$
  \cmark.}
\end{workedexample}

\begin{workedexample}{Landing in a moving cart}\label{ch1:we:cart}
A toy cart rolls across the floor, away from a bench, at a uniform velocity of
\qty{0.60}{m.s^{-1}}. A marble rolls horizontally off the edge of the bench, towards the cart.
At that moment the middle of the cart is \qty{0.50}{m} (horizontally) from the edge of the
bench. The top of the cart is \qty{0.80}{m} below the top of the bench. Find the launch speed
the marble needs to land in the middle of the cart.
\given{marble: $u_y = 0$, $s_y = \qty{-0.80}{m}$ (down to the top of the cart),
  $a_y = \qty{-9.8}{m.s^{-2}}$; \ cart: starts at $x_0 = \qty{0.50}{m}$ and moves away at
  \qty{0.60}{m.s^{-1}}. Positions $x$ are measured from the edge of the bench, forwards positive.}
\find{the marble's launch speed $u$}
\formula{marble, vertical (constant acceleration): $s_y = u_y t + \frac{1}{2}a_y t^2$; \
  marble, horizontal (uniform velocity): $x = ut$; \ cart (uniform velocity): $x = x_0 + 0.60\,t$}
\substitute{\textbf{Marble, vertical} (constant acceleration, $u_y = 0$): the time in the
  air.\par
  \eqlines{\qty{-0.80}{m} &= 0 + \tfrac{1}{2}(\qty{-9.8}{m.s^{-2}})\,t^2 = (\qty{-4.9}{m.s^{-2}})\,t^2\\
  t^2 &= \frac{\qty{-0.80}{m}}{\qty{-4.9}{m.s^{-2}}} = \qty{0.1633}{s^2}
    \quad\text{(divide both sides by $-4.9$)}\\
  t &= \sqrt{\qty{0.1633}{s^2}} = \qty{0.4041}{s} \quad\text{(square root)}}\par
  \textbf{Cart} (uniform velocity): where it is at this time.\par
  \eqlines{x_\mathrm{cart} &= \qty{0.50}{m} + (\qty{0.60}{m.s^{-1}})(\qty{0.4041}{s}) = \qty{0.7424}{m}}\par
  \textbf{Marble, horizontal} (uniform velocity): it must reach the same place at the same
  time, so $x = ut = x_\mathrm{cart}$.\par
  \eqlines{u &= \frac{x_\mathrm{cart}}{t} = \frac{\qty{0.7424}{m}}{\qty{0.4041}{s}} = \qty{1.837}{m.s^{-1}}}}
\answer{\final{$u = \qty{1.8}{m.s^{-1}}$} (2 significant figures)}
\checkstep{if the cart stood still at \qty{0.50}{m}, the marble would need only
  $0.50 \div 0.4041 = \qty{1.24}{m.s^{-1}}$. The cart moves away, so a faster launch makes sense
  \cmark. Same place at the same time: $1.837 \times 0.4041 = \qty{0.742}{m}$ \cmark.}
\end{workedexample}

\begin{mistakes}
\begin{itemize}
  \item \xmark\ Using $s = ut$ for a falling object. \cmark\ $s = ut$ needs $a = 0$. Vertically,
        use $s = ut + \frac{1}{2}at^2$.
  \item \xmark\ Starting two clocks at different moments. \cmark\ Use one $t$ for every object,
        counted from the same moment.
  \item \xmark\ ``A parcel dropped from a moving plane falls straight down, or falls behind.''
        \cmark\ It keeps the plane's horizontal velocity, so it stays under the plane (with no
        air resistance).
  \item \xmark\ Throwing away a positive root. \cmark\ Below the top, a ball passes each height
        twice. Pick the time the question asks for.
\end{itemize}
\end{mistakes}

\begin{examtip}
\begin{itemize}
  \item ``A ball is thrown straight up from a train that moves at a constant velocity. Describe
        its path as seen from the train and from the ground'' (3 marks). (1) Horizontally there
        is no force, so the ball keeps the train's velocity (uniform velocity). (2) From the
        train, it goes straight up and down and lands back in the hand. (3) From the ground,
        uniform velocity sideways plus constant acceleration downwards gives a parabola.
  \item In two-object calculations, write the model for each object, for example
        ``cart: $x = 0.50 + 0.60t$''. The marker can then follow your method. State your
        assumptions.
  \item Graph questions: the gradient of a velocity--time graph is the acceleration, and the
        area under it is the displacement.
\end{itemize}
\end{examtip}

\begin{quickcheck}
  \item A ball is thrown straight up at \qty{9.0}{m.s^{-1}}. Find (a) the time to reach the top,
        (b) the maximum height and (c) its velocity \qty{1.5}{s} after the throw.
  \item A cyclist rides at a uniform velocity of \qty{6.0}{m.s^{-1}}. She tosses a ball straight
        up (as she sees it) at \qty{3.5}{m.s^{-1}} and catches it at the same height. (a) How
        long is the ball in the air? (b) How far does she ride in that time? (c) Does the ball
        come down in front of her, behind her or in her hand? Explain.
  \item The velocity--time graph of a ball thrown straight up is a straight line. It goes from
        $+\qty{14.7}{m.s^{-1}}$ at $t = 0$ to $-\qty{14.7}{m.s^{-1}}$ at $t = \qty{3.0}{s}$. (a) Find
        the gradient, and say what it means. (b) Use areas to find the maximum height and the
        displacement at $t = \qty{3.0}{s}$.
  \item A marble leaves a bench horizontally at \qty{2.0}{m.s^{-1}}, \qty{1.25}{m} above the
        floor. A toy cart rolls \emph{towards} the bench at a uniform velocity of
        \qty{0.50}{m.s^{-1}}. Ignore the cart's height. How far from the bench must the cart be
        when the marble leaves, so that the marble lands in it?
\end{quickcheck}

%% ==================================================================
\section{Key experiment: testing the projectile model}\label{ch1:sec:experiments}
\covers{M5.IQ1.03}

\begin{keywords}
  \kw{Independent variable}{the one thing you change on purpose (here the launch speed $u$).}
  \kw{Dependent variable}{the thing you measure to see the effect (here the range $R$).}
  \kw{Controlled variables}{the things you keep the same, so the test is fair.}
  \kw{Random error}{an error that changes from trial to trial, sometimes too big and sometimes
    too small. Repeats and a mean reduce it.}
  \kw{Systematic error}{an error in the same direction every time. Repeats do not remove it.}
\end{keywords}

The syllabus asks you to collect \term{primary data}: data you measure yourself. You use it to
\term{validate} the relationships in this chapter. This means you test whether the predictions
match real measurements. For a horizontal launch from height $h$, the model predicts
\begin{equation}
  t = \sqrt{\frac{2h}{g}} \qquad\text{and}\qquad R = u\sqrt{\frac{2h}{g}}
  \label{ch1:eq:benchrange}
\end{equation}
(Not on the formula sheet: see the card ``Horizontal launch from a height'' in
\cref{ch1:sec:model}.) So a graph of $R$ against $u$ should be a straight line through the
origin. Its gradient should be $\sqrt{2h/g}$ (see the gradient toolkit in \cref{ch1:sec:model}).

\begin{experiment}{Range of a ball launched horizontally}
\aim{To test whether the projectile model predicts the range of a ball launched horizontally
  from a bench, for different launch speeds.}
\apparatus{steel ball (diameter about \qty{16}{mm}); vernier calipers (to measure the ball's
  diameter); curved ramp with a horizontal end; light gate and timer (a \term{light gate} times
  how long an object blocks a light beam); plumb line (a weighted string that hangs straight
  down); metre rule; spirit level (a tool with a bubble that shows when something is level);
  carbon paper on white paper (the ball leaves a dark dot where it lands); safety glasses
  (\cref{ch1:fig:apparatus}).
  \begin{figure}
  \centering
  \begin{tikzpicture}[x=0.8cm, y=0.8cm]
      \fill[ground] (-4.6,-0.22) rectangle (4.8,0);
      \draw[hscink, line width=0.6pt] (-4.6,0) -- (4.8,0);
      \fill[hscmuted!22] (-4.4,0) rectangle (-0.12,3.0);
      \draw[hscmuted, line width=0.6pt] (-4.4,0) rectangle (-0.12,3.0);
      \fill[hscmuted!40] (-3.8,3.0) -- (-3.8,5.0) .. controls (-2.8,3.2) and (-1.8,3.18)
        .. (0,3.18) -- (0,3.0) -- cycle;
      \draw[hscink, line width=1.2pt] (-3.8,5.0) .. controls (-2.8,3.2) and (-1.8,3.18) .. (0,3.18);
      \fill[object] (-3.22,4.42) circle (3.4pt);
      \fill[object] (0,3.31) circle (3.4pt);
      \draw[hscink, line width=1pt] (-0.62,3.18) -- (-0.62,3.85) -- (-0.22,3.85) -- (-0.22,3.18);
      \draw[hscforce, line width=0.6pt, dashed] (-0.62,3.5) -- (-0.22,3.5);
      \draw[guide, line width=0.8pt] (0,3.31) .. controls (1.3,3.31) and (2.4,2.0) .. (2.9,0.13);
      \fill[object] (2.9,0.13) circle (3.4pt);
      \draw[hscink, line width=0.4pt] (0.06,3.12) -- (0.06,0.3);
      \fill[hscink] (0.06,0.3) -- (-0.02,0.42) -- (0.14,0.42) -- cycle;
      \fill[hscink!85] (2.3,0) rectangle (3.6,0.07);
      \draw[guide] (0.2,3.18) -- (4.5,3.18);
      \draw[displacement, <->] (4.3,0) -- (4.3,3.18) node[midway, right, text=hscink, font=\small] {$h$};
      \draw[<->, line width=0.6pt, hscink] (0.06,-0.55) -- (2.9,-0.55)
        node[midway, below, font=\small] {$R$};
      \node[diagramlabel, font=\footnotesize, anchor=south west] at (-3.7,5.0) {ball released here};
      \node[diagramlabel, font=\footnotesize] at (-2.3,1.5) {bench};
      \node[diagramlabel, font=\footnotesize, anchor=south] (lg) at (0.3,4.35) {light gate (measures $u$)};
      \draw[hscmuted, line width=0.4pt] (lg.south) -- (-0.42,3.88);
      \node[diagramlabel, font=\footnotesize, anchor=west, align=left] (pl) at (0.3,1.0) {plumb\\line};
      \draw[hscmuted, line width=0.4pt] (pl.west) -- (0.1,1.0);
      \node[diagramlabel, font=\footnotesize, anchor=north west] (cp) at (3.15,-0.3) {carbon paper};
      \draw[hscmuted, line width=0.4pt] (cp.north west) -- (3.3,0.05);
    \end{tikzpicture}
  \caption{Side view. The plumb line marks the floor point below the launch
    point.}\label{ch1:fig:apparatus}
  \end{figure}}
\method{\textbf{Variables:} independent: the launch speed $u$; dependent: the range $R$;
  controlled: $h$, the ball, the ramp and the landing surface.
  \begin{enumerate}
  \item Clamp the ramp so its end is horizontal (check with the spirit level). Measure $h$ from
        the bottom of the ball at the end of the ramp to the floor. Mark the floor point under
        the plumb line.
  \item Measure the ball's diameter with the vernier calipers. Set the light-gate beam at the
        height of the ball's centre, as close to the end of the ramp as possible.
  \item Release the ball from a marked point on the ramp. The light gate times how long the
        ball blocks its beam ($\Delta t$). Then $u = \text{diameter} \div \Delta t$. Example:
        $\Delta t = \qty{13.4}{ms} = \qty{0.0134}{s}$ gives
        $u = 0.0160 \div 0.0134 = \qty{1.19}{m.s^{-1}}$.
  \item Measure $R$ from the floor mark to the centre of the dot on the carbon paper.
  \item Do 5 trials from the same release point and find the mean $R$. Then do the same for
        4 more release points (5 speeds in total).
  \item Safety: wear safety glasses. Put a box behind the landing point to catch the ball
        after it bounces, so no one slips on it.
\end{enumerate}}
\results{typical results for $h = \qty{0.900}{m}$ (\cref{ch1:fig:results}). The model predicts
  $\sqrt{2 \times 0.900/9.8} = \qty{0.429}{s}$, so $R_\mathrm{predicted} = 0.429\,u$. The
  gradient's unit is $\si{m} \div \si{m.s^{-1}} = \si{s}$: it is the fall time.
  \par\smallskip
  \begin{minipage}[c]{0.46\linewidth}\centering\small
  \begin{tabular}{@{}cccc@{}}
    \toprule
    $u$ & mean $R$ & spread & predicted $R$\\
    (\si{m.s^{-1}}) & (\si{m}) & (\si{m}) & (\si{m})\\
    \midrule
    1.19 & 0.498 & 0.015 & 0.510\\
    1.51 & 0.631 & 0.019 & 0.647\\
    1.80 & 0.752 & 0.015 & 0.771\\
    2.08 & 0.876 & 0.007 & 0.891\\
    2.42 & 1.016 & 0.015 & 1.037\\
    \bottomrule
  \end{tabular}\par\smallskip
  {\footnotesize spread = largest $-$ smallest of the 5 trials}
  \end{minipage}\hfill
  \begin{minipage}[c]{0.52\linewidth}\centering
  \begin{tikzpicture}
    \begin{axis}[hsc, width=\linewidth, height=0.55\linewidth,
        xmin=0, xmax=2.7, ymin=0, ymax=1.2,
        xlabel={$u$ (\si{m.s^{-1}})}, ylabel={$R$ (\si{m})},
        xtick={0,0.5,...,2.5}, ytick={0,0.2,...,1.2},
        xticklabel style={/pgf/number format/.cd, fixed, fixed zerofill, precision=1},
        yticklabel style={/pgf/number format/.cd, fixed, fixed zerofill, precision=1}]
      \addplot[hscseriesB, dashed, line width=1pt, domain=0:2.6, samples=2] {0.4286*x};
      \addplot[hscseriesA, line width=1pt, domain=0:2.6, samples=2] {0.4227*x - 0.0062};
      \addplot[hscseriesA, only marks, mark=*, mark size=1.8pt]
        coordinates {(1.19,0.498) (1.51,0.631) (1.80,0.752) (2.08,0.876) (2.42,1.016)};
      \node[anchor=east, font=\footnotesize, text=hscink, align=right] at (axis cs:1.6,0.98)
        {model:\\ gradient \qty{0.429}{s}};
      \draw[hscmuted, line width=0.4pt] (axis cs:1.62,0.95) -- (axis cs:2.05,0.89);
      \node[anchor=west, font=\footnotesize, text=hscink, align=left] at (axis cs:1.45,0.28)
        {best fit:\\ gradient \qty{0.423}{s}};
      \draw[hscmuted, line width=0.4pt] (axis cs:1.55,0.40) -- (axis cs:1.35,0.56);
    \end{axis}
  \end{tikzpicture}
  \captionof{figure}{$R$ against $u$ (typical results).}\label{ch1:fig:results}
  \end{minipage}}
\conclusion{the points lie on a straight line very close to the origin. So $R \propto u$: $R$
  is \term{proportional} to $u$ (double $u$ and $R$ doubles), as \cref{ch1:eq:benchrange}
  predicts.\par
  The gradient, \qty{0.423}{s}, differs from the predicted \qty{0.429}{s} by
  $(0.429 - 0.423) \div 0.429 \times \qty{100}{\percent} = \qty{1.4}{\percent}$.\par
  To find $g$, square both sides of gradient $= \sqrt{2h/g}$: gradient$^2 = 2h/g$. Then
  multiply both sides by $g$ and divide both sides by gradient$^2$:
  \[ g = \frac{2h}{\text{gradient}^2} = \frac{2 \times 0.900}{0.423^2} = \frac{1.800}{0.1789}
     = \qty{10.06}{m.s^{-2}} \]
  This differs from \qty{9.8}{m.s^{-2}} by
  $(10.06 - 9.8) \div 9.8 \times \qty{100}{\percent} = \qty{2.7}{\percent}$ (about
  \qty{3}{\percent}).\par
  The data support (validate) the model for a dense ball at low speed. The small, steady
  shortfall is a systematic error in the measurements, not a failure of the model.}
\significance{the fall time is the same for every speed, so $R$ grows in step with $u$: the
  horizontal and vertical motions are independent.}
\validity{\begin{itemize}[beginpenalty=10000]
  \item \textbf{Validity} (does it test what it claims?): change only $u$; keep the ball, $h$
        and the floor the same. The ramp end is horizontal, so $u_y = 0$, as the model
        assumes. Use a dense ball, so air resistance is small, as the model assumes. Measure
        $u$ directly with the light gate. Do not calculate it from the ramp height with
        energy: a rolling ball also spins, so some of the energy goes into spinning. The real
        $u$ is then smaller than that calculation gives.
  \item \textbf{Reliability} (do repeats agree?): 5 trials per speed, small spread
        (\qty{0.019}{m} or less), then the mean.
  \item \textbf{Accuracy} (how close to the true value?): every $R$ is about \qty{2}{\percent}
        below the prediction. An error in the same direction every time is a systematic
        error. Air resistance is too small to cause it: for this steel ball it shortens $R$ by
        only about \qty{0.1}{\percent}. Likely causes: (1) The light-gate beam misses the centre
        of the ball. The ball then blocks the beam for less than its full diameter, so $u$
        comes out too big. A beam \qty{1.5}{mm} off centre makes $u$ about \qty{2}{\percent}
        too big. (2) $h$ is measured too big. (3) The end of the ramp tilts slightly down.
\end{itemize}}
\end{experiment}

Other ways to validate the model:
\begin{enumerate}
  \item \textbf{Range against launch angle.} Fire a ball from a spring launcher onto a landing
        surface at the same height as the launch point. Keep the spring setting the same, so
        $u$ stays the same. Change $\theta$ from \ang{10} to \ang{80} in \ang{10} steps, with 5
        trials each. Plot $R$ against $\sin 2\theta$. The model predicts a straight line
        through the origin with gradient $u^2/g$. The largest $R$ should be at \ang{45}, and
        \ang{30} and \ang{60} should give equal ranges. The launch and landing heights must
        match: if the ball lands below the launch point, the best angle is less than \ang{45}.
  \item \textbf{Range against launch height.} Use the bench set-up with one release point (one
        $u$) and 4 different bench heights. Plot $R$ against $\sqrt{h}$. The model predicts a
        straight line through the origin.
  \item \textbf{Video analysis.} Film a thrown ball and study the video frame by frame. The
        $v_y$--$t$ graph should be a straight line with gradient \qty{-9.8}{m.s^{-2}}, and the
        $v_x$--$t$ graph should be flat.
\end{enumerate}

\begin{examtip}
\textbf{Validity}: one variable changed, the others controlled, the model's assumptions met.
\textbf{Reliability}: repeated trials agree (repeat at least 3 times). \textbf{Accuracy}: close
to the accepted value (give a percentage difference). Never write only ``human error''. Name
the error, and say how your improvement reduces it.
\end{examtip}

\begin{quickcheck}
  \item In the spring-launcher investigation (range against launch angle), name the
        independent variable, the dependent variable and two controlled variables.
  \item A student finds $u$ from the release height using $\frac{1}{2}mv^2 = mg\Delta h$
        instead of a light gate. Explain why this makes the experiment less valid.
  \item Every measured range was about \qty{2}{\percent} below the prediction. Is this a random
        error or a systematic error? Suggest one cause.
\end{quickcheck}

%% ==================================================================
% Chapter summary (added by the whole-book consistency edit: "How to use this guide" says
% that every chapter ends with a short summary). Every statement is taught above.
\begin{summarybox}[unbreakable]{\Cref{ch1:chap} in five lines}
\begin{itemize}
  \item Split the launch velocity: $u_x = u\cos\theta$ and $u_y = u\sin\theta$. The horizontal
        and vertical motions share only the time $t$.
  \item Horizontally there is no force: uniform velocity, $s_x = u_x t$. Vertically
        $a_y = \qty{-9.8}{m.s^{-2}}$ all the time: use the three equations of motion.
  \item At the top only $v_y = 0$. On level ground, $h_\mathrm{max} = \frac{u_y^2}{2g}$,
        $t_\mathrm{flight} = \frac{2u_y}{g}$ and $R = \frac{u^2\sin 2\theta}{g}$, which is largest
        at \ang{45}.
  \item From a height, solve $-h = u_y t - \tfrac{1}{2}gt^2$ for $t$ and keep the positive
        root. A final velocity needs a size \emph{and} a direction.
  \item Air resistance makes the range shorter and the path lower and steeper. The model works
        best for dense, slow objects over short distances.
\end{itemize}
\end{summarybox}
